\chapter{Operators and Spectra}\label{ch:operators}

\begin{center}\itshape
It may well be the granddaddy of all the learn-by-doing advanced mathematics books.\\[2pt]
\upshape--- Stuart Jay Sidney, on Pólya and Szegő's \emph{Aufgaben und Lehrsätze aus der Analysis}
\end{center}

In 1915 the young George Pólya, then in Zürich, had a guess and no proof. Take a positive function $f$ on the circle, write its Fourier coefficients
in a square table whose entry in row $j$ and column $k$ is $\hat f(j-k)$, and compute the determinant of the $n\times n$ corner. Otto Toeplitz had studied
these tables a few years earlier. Pólya guessed that the $n$-th root of the determinant tends to the geometric mean of $f$, $\exp\int\log f$, and he did
something that mathematicians rarely do: he published the conjecture. ``I had no proof,'' he said much later, ``but I published the conjecture and the
young Szegő found the proof.''

Gábor Szegő was twenty, a student in Budapest who had spent semesters in Berlin and Göttingen. His proof appeared in the \emph{Mathematische Annalen}
in 1915. On 15 May of the same year he enlisted in the Austro-Hungarian cavalry rather than wait to be called up, and he stayed in uniform until the
spring of 1919, passing through the infantry, the artillery and the air force, where he met Richard von Mises and Theodore von Kármán. The University of
Vienna gave him his doctorate on 8 July 1918, for the paper of 1915.

After the war the guesser and the prover wrote a book together. Their \emph{Aufgaben und Lehrsätze aus der Analysis} (1925) is a sequence of problems so
well chosen that solving them in order amounts to discovering a large part of analysis. Pólya called it his best work, and Szegő's too. Problem books
like the one you are holding descend from it.

Toeplitz determinants had a second life. In 1948, at a conference at Cornell, and again in Florence in 1949, Lars Onsager announced a formula for the
spontaneous magnetisation of the two-dimensional Ising model, a toy magnet on a square grid. He did not publish a proof: the magnetisation was the limit of
a Toeplitz determinant, and the limit was not covered by Szegő's theorem, which only gives the determinant up to a factor $e^{o(n)}$. Bruria Kaufman
and Onsager spoke to Shizuo Kakutani, and Kakutani spoke to Szegő. In 1952 Szegő, by then at Stanford, found the missing factor for a very general class of
functions, the \emph{strong} Szegő limit theorem. That same year C.~N.~Yang published a derivation of the formula, ``a tour de force''.

This chapter is about the eigenvalues of operators that come from somewhere: Toeplitz matrices and matrices built from divisibility, the Laplacian of
a drum and of a grid, integral operators on an interval, and the covariance and transition operators of probability. The theme throughout is
\emph{counting}: how many eigenvalues lie in an interval, how fast they decay, what their product is, and what the count reveals.

\begin{center}
\fbox{\parbox{0.92\linewidth}{\textbf{Think first.} A real symmetric $1000\times1000$ matrix has entries $a_{jk}=1$ if $j=k$, $a_{jk}=\frac12$ if $|j-k|=1$, and
$0$ otherwise. Without computing anything, how many of its eigenvalues exceed $1$? How many exceed $\frac32$? (Problems~\ref{pb:tridiag} and~\ref{pb:szego}.)}}
\end{center}

\emph{Notation.} $e(x)=e^{2\pi ix}$. A \emph{real trigonometric polynomial} is $f(\theta)=\sum_{|k|\le d}c_ke(k\theta)$ with $c_{-k}=\overline{c_k}$, so that $f$ is real on $\R/\Z$;
its \emph{Toeplitz matrices} are the Hermitian matrices $T_n(f)=(c_{j-l})_{0\le j,l<n}$, with $c_k=0$ for $|k|>d$. The eigenvalues of a Hermitian matrix or of a compact
self-adjoint operator are listed with multiplicity, and for a positive compact operator $\lambda_1\ge\lambda_2\ge\dots\ge0$. All integrals over $\R/\Z$ are over one period.

\section{Toeplitz matrices}

\begin{exercise}[$\star$]\label{pb:tridiag}
\emph{The tridiagonal case.} Let $a\in\R$, $b>0$ and $f(\theta)=a+2b\cos2\pi\theta$.
\begin{enumerate}[label=(\alph*)]
\item Show that the eigenvalues of $T_n(f)$ are $a+2b\cos\frac{\pi j}{n+1}$, $1\le j\le n$, with eigenvectors $\bigl(\sin\frac{\pi jk}{n+1}\bigr)_{1\le k\le n}$.
\item Show that for every $x\in\R$, the proportion of eigenvalues of $T_n(f)$ that are at most $x$ tends to the measure of $\{\theta\in\R/\Z:f(\theta)\le x\}$.
\item Answer the think box.
\end{enumerate}
\end{exercise}

\begin{exercise}[$\star\star\star$]\label{pb:szego}
\emph{Szegő's first limit theorem.} Let $f$ be a real trigonometric polynomial of degree $d$, with $m=\min f$ and $M=\max f$, and let $\lambda_1^{(n)},\dots,\lambda_n^{(n)}$ be the
eigenvalues of $T_n(f)$.
\begin{enumerate}[label=(\alph*)]
\item Show that all $\lambda_j^{(n)}$ lie in $[m,M]$, and that $|\langle T_n(g)x,y\rangle|\le\|g\|_\infty\|x\|\,\|y\|$ for every trigonometric polynomial $g$.
\item Show that $T_n(f)T_n(g)-T_n(fg)$ has at most $2d$ non-zero rows, for every trigonometric polynomial $g$.
\item Deduce that $\bigl|\operatorname{tr}T_n(f)^k-n\int_0^1f^k\bigr|\le4(k-1)d\,\|f\|_\infty^k$ for all $k\ge1$ and $n\ge1$.
\item Deduce that $\frac1n\sum_{j=1}^nF(\lambda_j^{(n)})\to\int_0^1F(f(\theta))\dd\theta$ for every continuous $F$ on $[m,M]$.
\end{enumerate}
\end{exercise}

\begin{exercise}[$\star\star$]\label{pb:szegodet}
\emph{Szegő's theorem of 1915.} Let $f$ be a real trigonometric polynomial.
\begin{enumerate}[label=(\alph*)]
\item Show that if $f>0$ on $\R/\Z$, then $\bigl(\det T_n(f)\bigr)^{1/n}\to\exp\int_0^1\log f$.
\item Let $f(\theta)=|1-re(\theta)|^2$ with $0<r<1$. Show that $\det T_n(f)=\frac{1-r^{2n+2}}{1-r^2}$ and $\int_0^1\log f=0$. So $\det T_n(f)$ divided by $\exp\bigl(n\int\log f\bigr)$
tends to $\frac1{1-r^2}$, a constant that (a) cannot see.
\item Let $f(\theta)=|1-e(\theta)|^2$, which vanishes at $0$. Show that $\det T_n(f)=n+1$, and that the conclusion of (a) still holds.
\end{enumerate}
\end{exercise}

\begin{exercise}[$\star\star$]\label{pb:kms}
\emph{The Kac--Murdock--Szegő matrices.} Let $-1<\rho<1$ and $A_n=(\rho^{|j-k|})_{0\le j,k<n}$.
\begin{enumerate}[label=(\alph*)]
\item Show that $(1-\rho^2)A_n^{-1}$ is the tridiagonal matrix with diagonal $1,1+\rho^2,\dots,1+\rho^2,1$ and $-\rho$ next to the diagonal, and that
$\det A_n=(1-\rho^2)^{n-1}$.
\item Show that $A_n=T_n(P_\rho)$ for the \emph{Poisson kernel} $P_\rho(\theta)=\frac{1-\rho^2}{1-2\rho\cos2\pi\theta+\rho^2}$, and deduce that the eigenvalues of $A_n$ lie in
$\bigl[\frac{1-|\rho|}{1+|\rho|},\frac{1+|\rho|}{1-|\rho|}\bigr]$.
\item Check that $\int_0^1\log P_\rho=\log(1-\rho^2)$, in agreement with (a) and Problem~\ref{pb:szegodet}(a), although $P_\rho$ is not a trigonometric polynomial.
\end{enumerate}
\end{exercise}

\begin{figure}[ht]
\centering
\fig{ch03-szego}
\caption{Szegő's theorem at work for $f(\theta)=1+\cos2\pi\theta+\frac12\cos6\pi\theta$ and $n=400$ (Problem~\ref{pb:szego}): the sorted eigenvalues of $T_n(f)$ (dots, one in ten)
against the increasing rearrangement of $f$ (curve). The $j$-th eigenvalue is close to the value that $f$ exceeds on a proportion $1-\frac jn$ of the circle.}\label{fig:szego}
\end{figure}

\section{Matrices from divisibility and from grids}

\begin{exercise}[$\star\star$]\label{pb:redheffer}
\emph{Redheffer's matrix.} Let $R_n$ be the $n\times n$ matrix with $(R_n)_{ij}=1$ if $i\mid j$ or $j=1$, and $0$ otherwise. Let $f_k(j)$ be the number of ways of writing $j$ as an
ordered product of $k$ integers $\ge2$, and $K=\lfloor\log_2n\rfloor$.
\begin{enumerate}[label=(\alph*)]
\item Show that
\[
\det(xI-R_n)=(x-1)^n-\sum_{k=1}^K(x-1)^{n-k-1}S_k,\qquad S_k=\sum_{j=2}^nf_k(j).
\]
\item Deduce that $\det R_n=M(n)=\sum_{j\le n}\mu(j)$, the Mertens function.
\item Show that $1$ is an eigenvalue of $R_n$ of algebraic multiplicity exactly $n-K-1$.
\end{enumerate}
\end{exercise}

\begin{exercise}[$\star$]\label{pb:minmatrix}
\emph{A discrete Brownian motion.} Let $A_n=(\min(i,j))_{1\le i,j\le n}$, the covariance matrix of a simple random walk at times $1,\dots,n$.
\begin{enumerate}[label=(\alph*)]
\item Show that $A_n^{-1}$ is tridiagonal, with diagonal $2,\dots,2,1$ and $-1$ next to the diagonal.
\item Show that the eigenvalues of $A_n^{-1}$ are $4\sin^2\frac{(2k-1)\pi}{2(2n+1)}$, $1\le k\le n$, with eigenvectors $\bigl(\sin\frac{(2k-1)\pi j}{2n+1}\bigr)_{1\le j\le n}$.
\item Deduce the identity $\sum_{k=1}^n\frac1{4\sin^2\frac{(2k-1)\pi}{2(2n+1)}}=\frac{n(n+1)}2$, and from it $\sum_{k\ge1}\frac1{(2k-1)^2}=\frac{\pi^2}8$.
\end{enumerate}
\end{exercise}

\begin{exercise}[Journey]\label{pb:gridtrees}
\emph{Spanning trees of the grid.} Let $\Gamma_{m,n}$ be the $m\times n$ grid graph (vertices $\{0,\dots,m-1\}\times\{0,\dots,n-1\}$, edges between points at distance $1$), and
$\tau_{m,n}$ its number of spanning trees.
\begin{enumerate}[label=(\alph*)]
\item Show that the eigenvalues of the Laplacian of the path with $m$ vertices are $2-2\cos\frac{\pi j}m$, $0\le j<m$, and that those of $\Gamma_{m,n}$ are the sums
$\bigl(2-2\cos\frac{\pi j}m\bigr)+\bigl(2-2\cos\frac{\pi k}n\bigr)$.
\item Show that $\prod_{k=1}^{n-1}\bigl(2\cosh t-2\cos\frac{\pi k}n\bigr)=\frac{\sinh nt}{\sinh t}$ for $t>0$, and deduce, with the matrix-tree theorem,
\[
\tau_{m,n}=\prod_{j=1}^{m-1}\frac{\sinh nt_j}{\sinh t_j},\qquad\cosh t_j=2-\cos\frac{\pi j}m.
\]
Check that $\tau_{2,2}=4$ and $\tau_{3,3}=192$.
\item Show that $\frac{\log\tau_{m,n}}{mn}\to\frac1\pi\int_0^\pi\operatorname{arccosh}(2-\cos\theta)\dd\theta$ as $m,n\to\infty$.
\item Show that this limit is $\frac{4G}\pi=1.16624\ldots$, where $G=\sum_{k\ge0}\frac{(-1)^k}{(2k+1)^2}$ is Catalan's constant. So a large grid has about $3.2099^{mn}$ spanning trees.
\end{enumerate}
\end{exercise}

\section{Hearing the shape of a drum}

For a bounded open set $\Omega\subseteq\R^2$, the \emph{Dirichlet} eigenvalues are the $\lambda$ for which $-\Delta u=\lambda u$ has a non-zero solution vanishing on the boundary
(a drum held at its rim), and the \emph{Neumann} eigenvalues those with zero normal derivative (a free membrane). For a rectangle $]0,a[\times]0,b[$, separation of variables gives
them explicitly: the Dirichlet eigenfunctions are $\sin\frac{m\pi x}a\sin\frac{n\pi y}b$, $m,n\ge1$, the Neumann ones are $\cos\frac{m\pi x}a\cos\frac{n\pi y}b$, $m,n\ge0$, and in both
cases the eigenvalue is $\pi^2\bigl(\frac{m^2}{a^2}+\frac{n^2}{b^2}\bigr)$. We write $N_D(\lambda)$ and $N_N(\lambda)$ for the number of eigenvalues at most $\lambda$, with multiplicity.

\begin{exercise}[$\star$]\label{pb:squaredrum}
\emph{The square drum and sums of two squares.} Let $\Omega=]0,1[^2$.
\begin{enumerate}[label=(\alph*)]
\item Show that the functions $2\sin m\pi x\sin n\pi y$, $m,n\ge1$, form a Hilbert basis of $L^2(\Omega)$.
\item Show that the multiplicity of the Dirichlet eigenvalue $\pi^2k$ is $\sum_{d\mid k}\chi_4(d)-\mathbf 1_{k\text{ is a square}}$, where $\chi_4$ is the non-trivial character modulo $4$.
\item Deduce that the multiplicities are unbounded. Show that the smallest eigenvalues of multiplicity $3$ and $4$ are $50\pi^2$ and $65\pi^2$.
\end{enumerate}
\end{exercise}

\begin{exercise}[$\star\star$]\label{pb:weylrect}
\emph{Weyl's law for a rectangle.} Let $\Omega=]0,a[\times]0,b[$.
\begin{enumerate}[label=(\alph*)]
\item Show that $N_N(\lambda)-N_D(\lambda)=\bigl\lfloor\frac{a\sqrt\lambda}\pi\bigr\rfloor+\bigl\lfloor\frac{b\sqrt\lambda}\pi\bigr\rfloor+1$.
\item Show that $N_D(\lambda)\le\frac{ab}{4\pi}\lambda\le N_N(\lambda)$ for all $\lambda>0$.
\item Deduce that $N_D(\lambda)=\frac{ab}{4\pi}\lambda+O(\sqrt\lambda)$: either spectrum determines the area. Show that the two spectra together determine the perimeter.
\end{enumerate}
\end{exercise}

\begin{figure}[ht]
\centering
\fig{ch03-weyl}
\caption{The counting functions $N_D$ (blue) and $N_N$ (red) of the rectangle $]0,1[\times]0,\sqrt2[$ (Problem~\ref{pb:weylrect}), the Weyl term $\frac{ab}{4\pi}\lambda$ (black) squeezed between
them, and the two-term approximations $\frac{ab}{4\pi}\lambda\mp\frac{a+b}{2\pi}\sqrt\lambda$ (dashed), which the staircases follow on average (Problem~\ref{pb:heatrect}).}\label{fig:weyl}
\end{figure}

\begin{exercise}[$\star\star\star$]\label{pb:heatrect}
\emph{One can hear the shape of a rectangle.} Let $\Omega=]0,a[\times]0,b[$ and $Z_D(t)=\sum e^{-t\lambda}$, the sum over the Dirichlet eigenvalues with multiplicity.
\begin{enumerate}[label=(\alph*)]
\item Show that for $s>0$, $\sum_{m\ge1}e^{-\pi sm^2}=\frac1{2\sqrt s}-\frac12+\varepsilon(s)$ with $0\le\varepsilon(s)\le\frac{e^{-\pi/s}}{\sqrt s\,(1-e^{-\pi/s})}$.
\item Deduce that as $t\to0^+$,
\[
Z_D(t)=\frac{ab}{4\pi t}-\frac{a+b}{4\sqrt{\pi t}}+\frac14+O\Bigl(\frac{e^{-c/t}}t\Bigr),\qquad c=\min(a,b)^2.
\]
\item Deduce that the Dirichlet spectrum of a rectangle determines its sides. What does the Neumann spectrum give?
\end{enumerate}
\end{exercise}

\begin{exercise}[$\star\star$]\label{pb:phasespace}
\emph{One eigenvalue per area $2\pi$.} On the functions $p(x)e^{-x^2/2}$, $p$ a polynomial, let $a=x+\frac{\mathrm d}{\mathrm dx}$, $a^\dagger=x-\frac{\mathrm d}{\mathrm dx}$ and
$H=-\frac{\mathrm d^2}{\mathrm dx^2}+x^2$, the quantum harmonic oscillator.
\begin{enumerate}[label=(\alph*)]
\item Show that $H=a^\dagger a+1$ and $Ha^\dagger=a^\dagger(H+2)$. Deduce that $\varphi_n=(a^\dagger)^ne^{-x^2/2}$ satisfies $H\varphi_n=(2n+1)\varphi_n$, and that $\varphi_n=h_n(x)e^{-x^2/2}$ with $h_n$ a
polynomial of degree $n$.
\item Show that the $\varphi_n/\|\varphi_n\|_2$ form a Hilbert basis of $L^2(\R)$.
\item Let $N(\lambda)$ be the number of eigenvalues $2n+1\le\lambda$. Show that $N(\lambda)\sim\frac1{2\pi}\operatorname{area}\{(x,\xi)\in\R^2:\xi^2+x^2\le\lambda\}$, and that the same rule gives the
right asymptotics for $-u''$ on $[0,L]$ with $u(0)=u(L)=0$.
\item What does the rule predict for $-u''+x^4u$? (A numerical calculation tests the prediction.)
\end{enumerate}
\end{exercise}

\begin{exercise}[$\star$]\label{pb:nodal}
\emph{Nodal domains of the square.} The nodal domains of an eigenfunction $u$ on $\Omega=]0,1[^2$ are the connected components of $\{u\ne0\}$.
\begin{enumerate}[label=(\alph*)]
\item Show that every Dirichlet eigenfunction for the eigenvalue $5\pi^2$ has exactly two nodal domains.
\item Show that $\sin m\pi x\sin n\pi y$ has $mn$ nodal domains.
\item Courant's theorem says that an eigenfunction for $\lambda_k$, where $k$ is the first index of that eigenvalue in the ordered list, has at most $k$ nodal domains. Check it for the functions
of (b) with eigenvalue at most $20\pi^2$, and find when equality holds.
\end{enumerate}
\end{exercise}

\section{Integral operators}

\begin{exercise}[$\star\star$]\label{pb:eulerrational}
\emph{Euler's theorem from a string.} Let $G(x,y)=\min(x,y)-xy$ on $[0,1]^2$, $G_1=G$ and $G_{k+1}(x,y)=\int_0^1G_k(x,z)G(z,y)\dd z$.
\begin{enumerate}[label=(\alph*)]
\item Show that on each of the triangles $x\le y$ and $x\ge y$, every $G_k$ is a polynomial in $x,y$ with rational coefficients.
\item The operator with kernel $G$ has the eigenvalues $(n\pi)^{-2}$, $n\ge1$, with eigenfunctions $\sqrt2\sin n\pi x$ (Tome~7, Exercise~4.6). Show, using only Hilbert--Schmidt norms, that
$\int_0^1G_k(x,x)\dd x=\sum_{n\ge1}(n\pi)^{-2k}$ for $k\ge2$.
\item Deduce Euler's theorem: $\zeta(2k)/\pi^{2k}$ is rational for every $k\ge1$.
\item Show that $G_2(x,y)=\frac16x(1-y)(2y-x^2-y^2)$ for $x\le y$, and deduce $\zeta(6)=\frac{\pi^6}{945}$.
\end{enumerate}
\end{exercise}

\begin{exercise}[$\star\star\star$]\label{pb:eigendecay}
\emph{Smooth kernels have small eigenvalues.} Let $k$ be a continuous positive kernel on $[0,1]$ (Tome~7, Theorem~4.14) with $|k(x,y)-k(x',y')|\le L(|x-x'|+|y-y'|)$, and
$\lambda_1\ge\lambda_2\ge\dots$ the eigenvalues of its operator $K$.
\begin{enumerate}[label=(\alph*)]
\item Let $P$ be the orthogonal projection on a subspace of dimension $m$. Show that $\operatorname{tr}(PKP)\le\lambda_1+\dots+\lambda_m$.
\item Taking for $P$ the projection on the functions constant on each interval $[\frac{i-1}m,\frac im[$, show that $\sum_{j>m}\lambda_j\le\frac L{3m}$.
\item Deduce that $\lambda_n\le\frac{3L}{n^2}$ for $n\ge2$, and show with $k(x,y)=\min(x,y)$ that the exponent $2$ cannot be improved.
\end{enumerate}
\end{exercise}

\begin{exercise}[$\star\star\star$]\label{pb:slepian}
\emph{How many numbers describe a band-limited signal?} Let $c>0$ and $K_c$ be the operator on $L^2(]-1,1[)$ with kernel $\frac{\sin c(x-y)}{\pi(x-y)}$ (equal to $\frac c\pi$ on the diagonal),
with eigenvalues $\lambda_1\ge\lambda_2\ge\dots$
\begin{enumerate}[label=(\alph*)]
\item Show that $0\le\langle K_cf,f\rangle\le\|f\|_2^2$, so that all eigenvalues lie in $[0,1]$.
\item Show that $\sum_n\lambda_n=\frac{2c}\pi$.
\item Show that $\sum_n\lambda_n^2=\frac{2c}\pi-\frac{\log4c+\gamma+1}{\pi^2}+o(1)$ as $c\to\infty$.
\item Deduce that the number of eigenvalues larger than $\frac12$ is $\frac{2c}\pi+O(\log c)$, and that for fixed $\varepsilon\in\,]0,\frac12[$ only $O(\log c)$ eigenvalues lie in $[\varepsilon,1-\varepsilon]$.
\end{enumerate}
\end{exercise}

\begin{figure}[ht]
\centering
\fig{ch03-slepian}
\caption{The eigenvalues of $K_c$ for $c=20$ (blue) and $c=40$ (red), plotted against $n$ (Problem~\ref{pb:slepian}). They stay close to $1$, plunge near $n=\frac{2c}\pi$ (dashed lines,
$12.7$ and $25.5$), and are then tiny. The width of the plunge grows only like $\log c$.}\label{fig:slepian}
\end{figure}

\begin{exercise}[$\star\star\star$]\label{pb:jentzsch}
\emph{Jentzsch's theorem.} Let $k$ be real, continuous and symmetric on $[0,1]^2$, with $k>0$ everywhere, and $K$ its operator on real $L^2([0,1])$. Let $\lambda_1=\sup_{\|f\|=1}\langle Kf,f\rangle$.
\begin{enumerate}[label=(\alph*)]
\item Show that $\lambda_1>0$ is an eigenvalue, and that if $Kf=\lambda_1f$ then $K|f|=\lambda_1|f|$.
\item Show that every eigenfunction for $\lambda_1$ is continuous and either everywhere positive or everywhere negative, and that $\lambda_1$ is a simple eigenvalue.
\item Show that every other eigenvalue $\mu$ satisfies $|\mu|<\lambda_1$.
\item Show that $\min_x\frac{(Kg)(x)}{g(x)}\le\lambda_1\le\max_x\frac{(Kg)(x)}{g(x)}$ for every continuous $g>0$, and that $\lambda_1\ge\iint k$.
\end{enumerate}
\end{exercise}

\section{Spectra of chance}

\begin{exercise}[$\star\star\star$]\label{pb:prediction}
\emph{The Kolmogorov--Szegő formula.} Let $f$ be an even real trigonometric polynomial with $f>0$, with coefficients $c_k$. For $n\ge1$ let $(X_{-n},\dots,X_0)$ be a centred Gaussian
vector with $\mathbb E[X_jX_l]=c_{j-l}$, and let $\sigma_n^2$ be the minimum of $\mathbb E\bigl[(X_0-\sum_{k=1}^na_kX_{-k})^2\bigr]$ over $a\in\R^n$: the error of the best linear prediction of $X_0$ from the
last $n$ values. Put $\sigma_0^2=c_0$.
\begin{enumerate}[label=(\alph*)]
\item Show that the minimum is attained and that $\sigma_n^2=\frac{\det T_{n+1}(f)}{\det T_n(f)}$.
\item Show that $\sigma_n^2$ decreases and that $\sigma_n^2\to\exp\int_0^1\log f$.
\item Compute $\sigma_n^2$ when the covariances are $\rho^{|k|}$ (Problem~\ref{pb:kms}), and when $f(\theta)=|1-e(\theta)|^2$. Interpret.
\end{enumerate}
\end{exercise}

\begin{exercise}[$\star\star$]\label{pb:levinson}
\emph{Reflection coefficients.} Let $(\gamma(k))_{k\in\Z}$ be real with $\gamma(-k)=\gamma(k)$, such that every matrix $(\gamma(j-l))_{0\le j,l\le n}$ is positive definite. Let $a^{(n)}\in\R^n$ be the
coefficients of the best linear prediction of $X_0$ from $X_{-1},\dots,X_{-n}$, as in Problem~\ref{pb:prediction}, $\sigma_n^2$ its error, and
$\kappa_{n+1}=\sigma_n^{-2}\bigl(\gamma(n+1)-\sum_{k=1}^na_k^{(n)}\gamma(n+1-k)\bigr)$.
\begin{enumerate}[label=(\alph*)]
\item Show the Levinson--Durbin recursion: $a^{(n+1)}_{n+1}=\kappa_{n+1}$, $a^{(n+1)}_k=a^{(n)}_k-\kappa_{n+1}a^{(n)}_{n+1-k}$ for $1\le k\le n$, and $\sigma_{n+1}^2=(1-\kappa_{n+1}^2)\sigma_n^2$. Deduce
$|\kappa_n|<1$ and $\sigma_n^2=\gamma(0)\prod_{k\le n}(1-\kappa_k^2)$.
\item Conversely, show that for every $\gamma(0)>0$ and every sequence $(\kappa_n)_{n\ge1}$ in $]-1,1[$ there is exactly one such sequence $\gamma$.
\item Find the $\kappa_n$ for $\gamma(k)=\rho^{|k|}$, and for $\gamma(0)=2$, $\gamma(\pm1)=-1$, $\gamma(k)=0$ otherwise.
\end{enumerate}
\end{exercise}

\begin{exercise}[$\star\star$]\label{pb:quadratic}
\emph{Estimating a covariance.} Let $f,g$ be even real trigonometric polynomials with $f>0$, and $X=(X_1,\dots,X_n)$ a centred Gaussian vector with covariance matrix $T_n(f)$. Let
$Q_n=\frac1nX^{\mathsf T}T_n(g)X$.
\begin{enumerate}[label=(\alph*)]
\item Show that $\mathbb EQ_n=\frac1n\operatorname{tr}\bigl(T_n(g)T_n(f)\bigr)=\int_0^1fg+O(\frac1n)$.
\item Show that $\operatorname{Var}Q_n=\frac2{n^2}\operatorname{tr}\bigl((T_n(g)T_n(f))^2\bigr)\le\frac2n\|f\|_\infty^2\|g\|_\infty^2$.
\item Deduce that for each $k\ge0$, $\frac1n\sum_{j=1}^{n-k}X_jX_{j+k}\to c_k$ in $L^2$, where $c_k$ is the $k$-th coefficient of $f$.
\end{enumerate}
\end{exercise}

\begin{exercise}[$\star\star$]\label{pb:ruin}
\emph{How long does a gambler last?} A simple random walk starts at $x\in\{1,\dots,N-1\}$ and stops at the time $T$ when it first hits $0$ or $N$.
\begin{enumerate}[label=(\alph*)]
\item Show that $P_x(T>n)=\sum_{k=1}^{N-1}c_k(x)\cos^n\frac{\pi k}N$ with $c_k(x)=\frac2N\sin\frac{\pi kx}N\sum_{j=1}^{N-1}\sin\frac{\pi kj}N$.
\item Show that $P_x(T>2n)\sim C_x\cos^{2n}\frac\pi N$ with $C_x>0$. So the walk survives $n$ steps with probability about $e^{-\pi^2n/(2N^2)}$ when $N$ is large.
\item Deduce that $\sum_{k=1}^{N-1}\frac{c_k(x)}{1-\cos(\pi k/N)}=x(N-x)$.
\end{enumerate}
\end{exercise}

\begin{exercise}[$\star\star$]\label{pb:randomcirculant}
\emph{A random circulant matrix.} Let $\varepsilon_0,\dots,\varepsilon_{n-1}$ be independent and uniform in $\{-1,1\}$, and $C$ the circulant matrix with first row $(\varepsilon_0,\dots,\varepsilon_{n-1})$.
\begin{enumerate}[label=(\alph*)]
\item Show that $C$ is normal, with eigenvalues $\lambda_k=\sum_j\varepsilon_je(\frac{jk}n)$, $0\le k<n$, so that $\|C\|=\max_k|\lambda_k|$.
\item Show that $\|C\|\ge\sqrt n$ always.
\item Show that $P(\|C\|\ge t)\le4ne^{-t^2/(4n)}$, and deduce that $P\bigl(\|C\|\ge\sqrt{8n\log n}\bigr)\le\frac4n$.
\end{enumerate}
\end{exercise}

\begin{exercise}[$\star\star\star$]\label{pb:multitype}
\emph{Branching with several types.} Individuals have types $1,\dots,d$. Each individual of type $i$ has, independently of everything else, a random vector of children with
distribution depending only on $i$, mean $(M_{i1},\dots,M_{id})$ and finite second moments. Let $Z_n\in\N^d$ (a row vector) count the individuals of each type in generation $n$, with $Z_0$ fixed and
non-zero, and $\rho$ the spectral radius of $M$.
\begin{enumerate}[label=(\alph*)]
\item Show that $\mathbb E[Z_{n+1}\mid Z_n]=Z_nM$, and that if $\rho<1$ the population dies out almost surely.
\item Suppose $M_{ij}>0$ for all $i,j$ and $\rho>1$, and let $v>0$ be a Perron eigenvector, $Mv=\rho v$. Show that $W_n=\rho^{-n}Z_nv$ is a martingale bounded in $L^2$.
\item Deduce that the population survives forever with positive probability, and grows like $\rho^n$ on a set of positive probability.
\end{enumerate}
\end{exercise}

\section*{Hints}

\begin{hint}{pb:tridiag}
For (a), check the equation $bv_{k-1}+av_k+bv_{k+1}=\lambda v_k$ with $v_0=v_{n+1}=0$, using $\sin(\alpha-\beta)+\sin(\alpha+\beta)=2\sin\alpha\cos\beta$. For (b), count the $j$ with
$\cos\frac{\pi j}{n+1}\le\frac{x-a}{2b}$.
\end{hint}

\begin{hint}{pb:szego}
For (a), $\langle T_n(g)x,y\rangle=\int_0^1g(\theta)P_x(\theta)\overline{P_y(\theta)}\dd\theta$ with $P_x(\theta)=\sum_lx_le(l\theta)$. For (b), compare $\sum_{0\le k<n}$ with $\sum_{k\in\Z}$ in the product
of the matrices. For (c), telescope $T_n(f)^k-T_n(f^k)$ and use that a matrix of rank $r$ and norm $\le C$ has trace at most $rC$ in modulus. For (d), polynomials first, then Weierstrass.
\end{hint}

\begin{hint}{pb:szegodet}
For (a), $\frac1n\log\det T_n(f)=\frac1n\sum_j\log\lambda_j^{(n)}$. For (b), expand the determinant along the last row. For (c), the same recursion; and $\log|1-e(\theta)|^2$ is the limit of
$\log|1-re(\theta)|^2$.
\end{hint}

\begin{hint}{pb:kms}
For (a), multiply. For (b), sum the two geometric series $\sum_{k\ge0}\rho^ke(\pm k\theta)$, and use the Rayleigh quotient of Problem~\ref{pb:szego}(a), which only needs absolutely summable
coefficients. For (c), $\log P_\rho=\log(1-\rho^2)-\log|1-\rho e(\theta)|^2$.
\end{hint}

\begin{hint}{pb:redheffer}
Write $R_n=D+ue_1^{\mathsf T}$, where $D_{ij}=\mathbf 1_{i\mid j}$ and $u=(0,1,\dots,1)^{\mathsf T}$, and use $\det(A+ue_1^{\mathsf T})=\det A\,(1+e_1^{\mathsf T}A^{-1}u)$. The powers of $D-I$ count chains of divisors.
For (b), $\sum_k(-1)^kf_k$ is the Dirichlet inverse of $\mathbf 1$.
\end{hint}

\begin{hint}{pb:minmatrix}
For (a), multiply: $\min(i,j)$ is linear in $i$ on each side of $j$. For (b), as in Problem~\ref{pb:tridiag}, with the boundary condition at $n$ turned into $v_{n+1}=v_n$.
\end{hint}

\begin{hint}{pb:gridtrees}
For (b), both sides are polynomials in $\cosh t$ of degree $n-1$ with the same roots, the zeros of the Chebyshev polynomial $U_{n-1}$. For (c), $\log\frac{\sinh nt}{\sinh t}=nt+$ terms that
are small after division by $mn$. For (d), $2-\cos2u=1+2\sin^2u$ and $\operatorname{arccosh}(1+2s^2)=2\operatorname{arcsinh}s$; then differentiate $\int_0^{\pi/2}\operatorname{arcsinh}(\alpha\sin u)\dd u$ in $\alpha$.
\end{hint}

\begin{hint}{pb:squaredrum}
For (a), a function orthogonal to all products has, for almost every $x$, all its sine coefficients in $y$ equal to $0$ (Fubini). For (b), Jacobi's formula $r_2(k)=4\sum_{d\mid k}\chi_4(d)$
(Tome~1, Chapter~10).
\end{hint}

\begin{hint}{pb:weylrect}
For (b), attach to each lattice point $(m,n)$ in the quarter ellipse the unit square below and to its left (for $N_D$), or above and to its right (for $N_N$).
\end{hint}

\begin{hint}{pb:heatrect}
For (a), Poisson summation gives $\sum_{m\in\Z}e^{-\pi sm^2}=s^{-1/2}\sum_{m\in\Z}e^{-\pi m^2/s}$. For (b), $Z_D(t)$ is a product of two such sums. For (c), two numbers determine
$ab$ and $a+b$.
\end{hint}

\begin{hint}{pb:phasespace}
For (a), compute $a^\dagger au$, and use induction. For (b), rescale Tome~7, Theorem~5.5: the span of the $\varphi_n$ is the span of the $x^ke^{-x^2/2}$. For (d), compute the area of
$\{\xi^2+x^4\le\lambda\}$.
\end{hint}

\begin{hint}{pb:nodal}
For (a), $\sin2\alpha=2\sin\alpha\cos\alpha$ factors every eigenfunction as $2\sin\pi x\sin\pi y\,(\alpha\cos\pi y+\beta\cos\pi x)$.
\end{hint}

\begin{hint}{pb:eulerrational}
For (a), split $\int_0^1$ at $\min(x,y)$ and $\max(x,y)$. For (b), the Hilbert--Schmidt inner product of two operators with kernels is the integral of the product of the kernels, and
equals $\sum_n\langle Ae_n,Be_n\rangle$ in an orthonormal basis.
\end{hint}

\begin{hint}{pb:eigendecay}
For (a), the non-zero eigenvalues of $PKP$ are at most those of $K$, one by one (Courant--Fischer). For (b), $\operatorname{tr}K=\int_0^1k(x,x)\dd x$ by Mercer, and $\operatorname{tr}(PKP)$ is an
average of $k$ over the squares of side $\frac1m$ on the diagonal.
\end{hint}

\begin{hint}{pb:slepian}
For (a), Plancherel. For (b), Mercer. For (c), the Hilbert--Schmidt norm is $\iint K^2$; substitute $u=x-y$, and use $\int_0^X\frac{\sin^2v}v\dd v=\frac12(\log2X+\gamma)+o(1)$. For (d), compare
$\#\{\lambda_n>\frac12\}$ with $\sum\lambda_n$ using $\sum\lambda_n(1-\lambda_n)$.
\end{hint}

\begin{hint}{pb:jentzsch}
For (a), $\langle K|f|,|f|\rangle\ge|\langle Kf,f\rangle|$. For (b), $K|f|>0$ everywhere. For (c), if $Kg=\mu g$ with $|\mu|=\lambda_1$, look at $|g|$ and at the case of equality in
$|\int kg|\le\int k|g|$. For (d), pair with the positive eigenfunction.
\end{hint}

\begin{hint}{pb:prediction}
For (a), the normal equations, and the Schur complement in $T_{n+1}(f)$. For (b), the minimum over more variables is smaller, and Cesàro means of a convergent sequence converge to the same
limit.
\end{hint}

\begin{hint}{pb:levinson}
By stationarity, the best prediction of $X_{-n-1}$ from $X_{-n},\dots,X_{-1}$ has the coefficients of $a^{(n)}$ in reverse order. Correct the forward prediction by a multiple of the backward
error.
\end{hint}

\begin{hint}{pb:quadratic}
For (a), Problem~\ref{pb:szego}(b). For (b), write $X=\Sigma^{1/2}Z$ and diagonalise. For (c), choose $g(\theta)=2\cos2\pi k\theta$.
\end{hint}

\begin{hint}{pb:ruin}
$P_x(T>n)=(Q^n\mathbf 1)(x)$ for the matrix $Q$ of the walk killed at $0$ and $N$, which is $T_{N-1}(\cos2\pi\theta)$. For (b), group the terms $k$ and $N-k$.
\end{hint}

\begin{hint}{pb:randomcirculant}
For (b), Parseval: $\sum_k|\lambda_k|^2=n\sum_j\varepsilon_j^2$. For (c), Hoeffding's inequality for the real and imaginary parts of each $\lambda_k$, and a union bound.
\end{hint}

\begin{hint}{pb:multitype}
For (a), $\mathbb P(Z_n\ne0)\le\mathbb E|Z_n|_1$. For (b), the conditional variance of $Z_{n+1}v$ given $Z_n$ is linear in $Z_n$, and $\mathbb E|Z_n|_1=O(\rho^n)$.
\end{hint}

\section*{Solutions}

\begin{solution}{pb:tridiag}
\idea A tridiagonal Toeplitz matrix is a discrete second derivative, and sines that vanish at both ends diagonalise it, exactly as for a vibrating string.

\writeup (a) Let $v_k=\sin\frac{\pi jk}{n+1}$ for $0\le k\le n+1$, so that $v_0=v_{n+1}=0$. Since $\sin(\alpha-\beta)+\sin(\alpha+\beta)=2\sin\alpha\cos\beta$,
$bv_{k-1}+av_k+bv_{k+1}=\bigl(a+2b\cos\frac{\pi j}{n+1}\bigr)v_k$ for $1\le k\le n$, which is $T_n(f)v=\lambda_jv$ because $f$ has coefficients $c_0=a$, $c_{\pm1}=b$. The $n$ values $\lambda_j$ are
distinct, since $\cos$ is injective on $]0,\pi[$, so they are all the eigenvalues.

(b) Let $y=\frac{x-a}{2b}$. If $|y|\le1$, then $\lambda_j\le x$ if and only if $j\ge\frac{n+1}\pi\arccos y$, which holds for $n(1-\frac1\pi\arccos y)+O(1)$ values of $j$. On the other side,
$\cos2\pi\theta>y$ exactly on an arc of length $\frac1\pi\arccos y$ around $0$, so $\{f\le x\}$ has measure $1-\frac1\pi\arccos y$. If $|y|>1$, both sides are $0$ or $1$.

(c) Here $f(\theta)=1+\cos2\pi\theta$, so $a=1$, $b=\frac12$, $n=1000$ and $\lambda_j=1+\cos\frac{\pi j}{1001}$. Then $\lambda_j>1$ if and only if $j<500.5$: $500$ eigenvalues. And
$\lambda_j>\frac32$ if and only if $\frac{\pi j}{1001}<\frac\pi3$, that is $j\le333$: $333$ eigenvalues, against $\frac{1000}3$ predicted by (b), since $f>\frac32$ on a third of the circle.

\tools{Tome~0, Chapter~7 and Tome~2, Chapter~3 (eigenvalues of symmetric matrices).}
\leads The same sines diagonalise the discrete Laplacian of a path with fixed ends, and they are the reason why the discrete sine transform, a cousin of the FFT of Tome~4, Chapter~7, solves
Poisson's equation on a grid in $O(n\log n)$ operations.
\end{solution}

\begin{solution}{pb:szego}
\idea A Toeplitz matrix is the compression of the multiplication by $f$ to the first $n$ Fourier modes. Compression and multiplication almost commute: powers of $T_n(f)$ differ from $T_n(f^k)$
only near the corners, so the traces agree up to $O(1)$, and traces of powers determine the distribution of the eigenvalues.

\writeup (a) For $x,y\in\C^n$ let $P_x(\theta)=\sum_lx_le(l\theta)$. Then, with $c_k(g)=\int_0^1ge(-k\theta)$,
\[
\int_0^1gP_x\overline{P_y}=\sum_{j,l}x_l\overline{y_j}\int_0^1ge((l-j)\theta)\dd\theta=\sum_{j,l}c_{j-l}(g)x_l\overline{y_j}=\langle T_n(g)x,y\rangle.
\]
By Cauchy--Schwarz and Parseval, $|\langle T_n(g)x,y\rangle|\le\|g\|_\infty\|P_x\|_2\|P_y\|_2=\|g\|_\infty\|x\|\,\|y\|$. For $g=f$ and $y=x$, $\langle T_n(f)x,x\rangle=\int f|P_x|^2$ lies between $m\|x\|^2$ and
$M\|x\|^2$, so the eigenvalues of the Hermitian matrix $T_n(f)$ lie in $[m,M]$.

(b) Write $\hat f(k)=c_k(f)$. Since $c_{j-l}(fg)=\sum_{k\in\Z}\hat f(j-k)\hat g(k-l)$,
\[
\bigl(T_n(f)T_n(g)-T_n(fg)\bigr)_{jl}=-\sum_{k<0\text{ or }k\ge n}\hat f(j-k)\hat g(k-l).
\]
A term is non-zero only if $|j-k|\le d$. For $k<0$ this forces $j<d$, and for $k\ge n$ it forces $j\ge n-d$. So only the first $d$ and the last $d$ rows can be non-zero.

(c) Telescoping,
\[
T_n(f)^k-T_n(f^k)=\sum_{i=1}^{k-1}T_n(f)^{k-1-i}\bigl(T_n(f)T_n(f^i)-T_n(f^{i+1})\bigr).
\]
By (b) each term has rank at most $2d$, and by (a) its norm is at most $\|f\|_\infty^{k-1-i}\cdot2\|f\|_\infty^{i+1}=2\|f\|_\infty^k$. A matrix $A$ of rank $r$ has $|\operatorname{tr}A|\le\sum_is_i(A)\le r\|A\|$,
where the $s_i$ are its singular values. So $|\operatorname{tr}T_n(f)^k-\operatorname{tr}T_n(f^k)|\le4(k-1)d\|f\|_\infty^k$, and $\operatorname{tr}T_n(f^k)=nc_0(f^k)=n\int_0^1f^k$.

(d) By (c), $\frac1n\sum_j(\lambda_j^{(n)})^k=\frac1n\operatorname{tr}T_n(f)^k\to\int f^k$ for each $k$, hence $\frac1n\sum_jp(\lambda_j^{(n)})\to\int p\circ f$ for every polynomial $p$. Given $F$ continuous on
$[m,M]$ and $\varepsilon>0$, choose a polynomial $p$ with $|F-p|\le\varepsilon$ on $[m,M]$ (Weierstrass, Tome~6, Corollary~4.20). Since the $\lambda_j^{(n)}$ and the values of $f$ lie in $[m,M]$, both
sides move by at most $\varepsilon$ when $p$ replaces $F$, so the $\limsup$ of the difference is at most $2\varepsilon$.

\tools{Tome~7, Chapter~2 (Parseval); Tome~2, Chapter~6 (singular values); Tome~6, Corollary~4.20 (Weierstrass).}
\leads Szegő proved (d) in 1920 for continuous $f$; the version for $f\in L^\infty$ and the extension to non-Hermitian matrices (where it fails in general: the eigenvalues of $T_n(e(\theta))$ are all $0$)
are surveyed in Böttcher and Silbermann, \emph{Introduction to Large Truncated Toeplitz Matrices} (1999). The trace argument is the moment method of probability, and it reappears in random matrix
theory.
\end{solution}

\begin{solution}{pb:szegodet}
\idea The logarithm of a determinant is the sum of the logarithms of the eigenvalues, and Szegő's first limit theorem averages exactly such sums. The examples show what the limit theorem cannot
see: a bounded factor, which is the subject of the strong theorem.

\writeup (a) The eigenvalues lie in $[m,M]$ with $m=\min f>0$ (Problem~\ref{pb:szego}(a)), and $\log$ is continuous there. So $\frac1n\log\det T_n(f)=\frac1n\sum_j\log\lambda_j^{(n)}\to\int\log f$ by
Problem~\ref{pb:szego}(d).

(b) $f=1+r^2-r\bigl(e(\theta)+e(-\theta)\bigr)$, so $T_n(f)$ is tridiagonal with $1+r^2$ on the diagonal and $-r$ next to it. Expanding along the last row,
$D_n=(1+r^2)D_{n-1}-r^2D_{n-2}$ with $D_0=1$, $D_1=1+r^2$. The characteristic roots are $1$ and $r^2$, and $D_n=\frac{1-r^{2n+2}}{1-r^2}$ satisfies the recursion and the initial values. Next,
$\log|1-re(\theta)|^2=-2\operatorname{Re}\sum_{k\ge1}\frac{r^k}ke(k\theta)$, uniformly convergent, so its integral is $0$. Hence $D_n/e^{n\int\log f}=D_n\to\frac1{1-r^2}$.

(c) Now $D_n=2D_{n-1}-D_{n-2}$ with $D_0=1$, $D_1=2$, so $D_n=n+1$ and $\frac1n\log D_n\to0$. For $\frac12\le r<1$, $|1-re(\theta)|^2=(1-r)^2+4r\sin^2\pi\theta$ lies between $2\sin^2\pi\theta$ and $4$, so
$|\log|1-re(\theta)|^2|\le\log4+|\log(2\sin^2\pi\theta)|$, which is integrable. By dominated convergence $\int\log|1-e(\theta)|^2=\lim_{r\to1}\int\log|1-re(\theta)|^2=0$, and the conclusion of (a)
holds, although $\min f=0$.

\tools{Problem~\ref{pb:szego}; Tome~6, Theorem~7.9 (dominated convergence).}
\leads (a) is Szegő's theorem of 1915, Pólya's conjecture. In (b), $\log f=\sum_kg_ke(k\theta)$ with $g_{\pm k}=-\frac{r^k}k$, and $\frac1{1-r^2}=\exp\sum_{k\ge1}kg_kg_{-k}$: this is the \emph{strong} Szegő limit
theorem of 1952, $\det T_n(f)\sim\exp\bigl(ng_0+\sum_{k\ge1}kg_kg_{-k}\bigr)$, in the simplest example. Symbols with zeros, like (c), are governed by the Fisher--Hartwig asymptotics, where $n+1$ is a
power of $n$ times a constant.
\end{solution}

\begin{solution}{pb:kms}
\idea The matrix $(\rho^{|j-k|})$ is the covariance of a stationary autoregressive sequence $X_k=\rho X_{k-1}+\xi_k$; its inverse only couples neighbours, because $X_k$ depends on the past only
through $X_{k-1}$.

\writeup (a) Let $B$ be the tridiagonal matrix of the statement. For an interior column $l$ and $j<l$, with $d=l-j$,
$(A_nB)_{jl}=\rho^d(1+\rho^2)-\rho\cdot\rho^{d-1}-\rho\cdot\rho^{d+1}=0$; the same for $j>l$; and $(A_nB)_{ll}=1+\rho^2-2\rho^2=1-\rho^2$. In the first column, $(A_nB)_{j0}=\rho^j-\rho\rho^{j-1}=0$ for $j\ge1$ and
$1-\rho^2$ for $j=0$; the last column is symmetric. So $A_nB=(1-\rho^2)I$. For the determinant, subtract $\rho$ times row $j-1$ from row $j$, for $j=n-1,\dots,1$: row $j$ becomes $0$ before
the diagonal and $(1-\rho^2)\rho^{k-j}$ in columns $k\ge j$. The matrix is now triangular with diagonal $1,1-\rho^2,\dots,1-\rho^2$.

(b) $\sum_{k\in\Z}\rho^{|k|}e(k\theta)=1+2\operatorname{Re}\frac{\rho e(\theta)}{1-\rho e(\theta)}=\frac{1-\rho^2}{|1-\rho e(\theta)|^2}=P_\rho(\theta)$, so the coefficients of $P_\rho$ are $\rho^{|k|}$. The computation of
Problem~\ref{pb:szego}(a) only uses $\sum_k|c_k|<\infty$ (to integrate term by term), so $\langle A_nx,x\rangle=\int P_\rho|P_x|^2$ lies between $\min P_\rho$ and $\max P_\rho$ times $\|x\|^2$, and these
are $\frac{1-\rho^2}{(1+|\rho|)^2}=\frac{1-|\rho|}{1+|\rho|}$ and $\frac{1+|\rho|}{1-|\rho|}$.

(c) $\int\log P_\rho=\log(1-\rho^2)-\int\log|1-\rho e(\theta)|^2=\log(1-\rho^2)$ by Problem~\ref{pb:szegodet}(b) (also for $\rho<0$, replacing $\theta$ by $\theta+\frac12$). And
$\frac1n\log\det A_n=\frac{n-1}n\log(1-\rho^2)\to\log(1-\rho^2)$.

\tools{Problems~\ref{pb:szego} and~\ref{pb:szegodet}; Tome~0, Chapter~7 (determinants).}
\leads The name comes from Kac, Murdock and Szegő, \emph{On the eigenvalues of certain Hermitian forms} (1953), who proved continuous analogues of Szegő's theorem for integral operators such as
$e^{-|x-y|}$ on $[0,T]$, the covariance of the stationary Ornstein--Uhlenbeck process. Problem~\ref{pb:slepian} is another case.
\end{solution}

\begin{solution}{pb:redheffer}
\idea Redheffer's matrix is the divisibility matrix plus a column of ones. The divisibility matrix is unipotent, the powers of its nilpotent part count chains of divisors, and the alternating sum of
those counts is the Möbius function.

\writeup (a) Let $D_{ij}=\mathbf 1_{i\mid j}$, an upper triangular matrix with ones on the diagonal, $N=D-I$, and $u=(0,1,\dots,1)^{\mathsf T}$. Since $(R_n)_{i1}=1$ for every $i$ but $D_{i1}=\mathbf 1_{i=1}$,
$R_n=D+ue_1^{\mathsf T}$. For $x\ne1$, $A=xI-D=(x-1)I-N$ is invertible with $A^{-1}=\sum_{k\ge0}\frac{N^k}{(x-1)^{k+1}}$, a finite sum as $N$ is nilpotent, and $\det A=(x-1)^n$. The matrix
determinant lemma gives $\det(A-ue_1^{\mathsf T})=\det A\,(1-e_1^{\mathsf T}A^{-1}u)$. Now $(N^k)_{1j}$ counts the chains $1=d_0\mid d_1\mid\dots\mid d_k=j$ with $d_{i-1}<d_i$, that is, the ordered
factorisations $j=\frac{d_1}{d_0}\cdots\frac{d_k}{d_{k-1}}$ into $k$ factors $\ge2$: $(N^k)_{1j}=f_k(j)$. So $e_1^{\mathsf T}N^ku=S_k$ for $k\ge1$, $e_1^{\mathsf T}u=0$, and
\[
\det(xI-R_n)=(x-1)^n\Bigl(1-\sum_{k\ge1}\frac{S_k}{(x-1)^{k+1}}\Bigr).
\]
Since $f_k(j)=0$ when $2^k>j$, only $k\le K$ occur. Both sides are polynomials that agree for $x\ne1$, hence everywhere.

(b) At $x=0$: $(-1)^n\det R_n=(-1)^n-\sum_k(-1)^{n-k-1}S_k$, so $\det R_n=1+\sum_{k\ge1}(-1)^kS_k=1+\sum_{j=2}^ng(j)$ with $g=\sum_{k\ge0}(-1)^kh^{*k}$, where $h=\mathbf 1-\delta$ ($\delta$ the unit
of Dirichlet convolution) and $f_k=h^{*k}$; for each $j$ the sum is finite. Then $g*\mathbf 1=g*(\delta+h)=\sum_k(-1)^kh^{*k}+\sum_k(-1)^kh^{*(k+1)}=\delta$, so $g=\mu$ (Tome~1, Chapter~4), and
$\det R_n=1+\sum_{j=2}^n\mu(j)=M(n)$.

(c) $\det(xI-R_n)=(x-1)^{n-K-1}q(x)$ with $q(x)=(x-1)^{K+1}-\sum_{k=1}^K(x-1)^{K-k}S_k$, and $q(1)=-S_K\le-f_K(2^K)=-1$ since $2^K\le n$. So the multiplicity is exactly $n-K-1$.

\tools{Tome~1, Chapter~4 (Dirichlet convolution, Möbius inversion); Tome~0, Chapter~7 (determinants).}
\leads Redheffer introduced the matrix in 1977. The Riemann hypothesis is equivalent to $M(n)=O(n^{1/2+\varepsilon})$ for every $\varepsilon>0$ (Littlewood, 1912), so it is a statement about the
determinant of a $0$--$1$ matrix. The other $K+1$ eigenvalues are the roots of $q$: Barrett, Forcade and Pollington (1988) and Vaughan (1993) showed that the spectral radius is of order $\sqrt n$,
the remaining ones being bounded, and numerical evaluation shows them.
\end{solution}

\begin{solution}{pb:minmatrix}
\idea $\min(i,j)$ is piecewise linear in $i$, so the second difference kills it except on the diagonal. The inverse is a discrete Laplacian with a fixed end at $0$ and a free end at $n$, and its
eigenvectors are sines that satisfy both boundary conditions.

\writeup (a) Let $B$ be the tridiagonal matrix of the statement. For $l<n$, $(A_nB)_{il}=2\min(i,l)-\min(i,l-1)-\min(i,l+1)$ (with $\min(i,0)=0$), the second difference at $l$ of $k\mapsto\min(i,k)$,
which is linear on $[0,i]$ and on $[i,n]$: it is $0$ unless $l=i$, where it is $2i-(i-1)-i=1$. For $l=n$, $(A_nB)_{in}=\min(i,n)-\min(i,n-1)=\mathbf 1_{i=n}$.

(b) Let $\alpha=\frac{(2k-1)\pi}{2n+1}$ and $v_j=\sin j\alpha$. For $1\le j<n$, $-v_{j-1}+2v_j-v_{j+1}=(2-2\cos\alpha)v_j=4\sin^2\frac\alpha2\,v_j$, with $v_0=0$. The last row asks for
$-v_{n-1}+v_n=4\sin^2\frac\alpha2\,v_n$, which by the same identity at $j=n$ means $v_{n+1}=v_n$, that is $\sin(n+1)\alpha=\sin n\alpha$; this holds because $(n+1)\alpha+n\alpha=(2k-1)\pi$. The $n$ values of
$\alpha$ lie in $]0,\pi[$ and give distinct eigenvalues.

(c) The eigenvalues of $A_n$ are the inverses, and $\operatorname{tr}A_n=1+2+\dots+n$. Dividing by $n^2$, the $k$-th term $\bigl(4n^2\sin^2\frac{(2k-1)\pi}{4n+2}\bigr)^{-1}$ tends to $\frac4{(2k-1)^2\pi^2}$, and by
$\sin x\ge\frac{2x}\pi$ on $[0,\frac\pi2]$ it is at most $\frac{(4n+2)^2}{16n^2(2k-1)^2}\le\frac9{4(2k-1)^2}$. Dominated convergence for series gives $\sum_k\frac4{(2k-1)^2\pi^2}=\frac12$.

\tools{Tome~0, Chapter~7 (inverse matrices); Tome~6, Theorem~7.9 (dominated convergence).}
\leads $A_n/n$ is the covariance of a random walk observed at times $\frac jn$, and $\frac1nA_n$ acting on vectors with the weight $\frac1n$ approximates the operator with kernel $\min(s,t)$: the eigenvalues
$\frac1{n^2}\bigl(4\sin^2\frac{(2k-1)\pi}{4n+2}\bigr)^{-1}\to\frac1{(k-\frac12)^2\pi^2}$ are those of Brownian motion (Tome~7, Example~4.15), and (c) is the discrete Karhunen--Loève expansion.
\end{solution}

\begin{solution}{pb:gridtrees}
\idea Kirchhoff's theorem turns trees into eigenvalues, the grid's eigenvalues are sums of those of two paths, and a Chebyshev product collapses one of the two products. What is left is a
Riemann sum, whose limit is an integral that a parameter trick evaluates.

\writeup (a) For the path, let $w(x)=\cos\frac{\pi j(x+1/2)}m$, $0\le x<m$. At interior points $2w(x)-w(x-1)-w(x+1)=(2-2\cos\frac{\pi j}m)w(x)$; at $x=0$ the Laplacian gives $w(0)-w(1)$, which is the
same expression because $w(-1)=w(0)$, and similarly $w(m)=w(m-1)$. The $m$ eigenvalues are distinct. The grid is the Cartesian product of two paths, so its Laplacian is
$L_m\otimes I+I\otimes L_n$, with eigenvectors $w_j\otimes w'_k$ and eigenvalues the sums.

(b) The Chebyshev polynomial $U_{n-1}$, defined by $U_{n-1}(\cos\varphi)=\frac{\sin n\varphi}{\sin\varphi}$, has degree $n-1$, leading coefficient $2^{n-1}$ (from $U_k=2yU_{k-1}-U_{k-2}$) and roots
$\cos\frac{\pi k}n$, $1\le k<n$. So $\prod_k(2y-2\cos\frac{\pi k}n)=U_{n-1}(y)$, and at $y=\cosh t=\cos(it)$ this is $\frac{\sinh nt}{\sinh t}$. Letting $t\to0$, $\prod_{k=1}^{n-1}(2-2\cos\frac{\pi k}n)=n$.
By the matrix-tree theorem (Tome~11, Corollary~A.5), $\tau_{m,n}$ is $\frac1{mn}$ times the product of the non-zero eigenvalues. For $j\ge1$, with $2\cosh t_j=4-2\cos\frac{\pi j}m$, the eigenvalues
are $2\cosh t_j-2\cos\frac{\pi k}n$, and their product over $0\le k<n$ is $(2\cosh t_j-2)\frac{\sinh nt_j}{\sinh t_j}$, where $2\cosh t_j-2=2-2\cos\frac{\pi j}m$. For $j=0$ the non-zero eigenvalues multiply
to $n$. Since $\prod_{j=1}^{m-1}(2-2\cos\frac{\pi j}m)=m$, $\tau_{m,n}=\frac{mn}{mn}\prod_j\frac{\sinh nt_j}{\sinh t_j}$. For $m=n=2$, $\cosh t_1=2$ and $\frac{\sinh2t}{\sinh t}=2\cosh t=4$. For $m=n=3$,
$\frac{\sinh3t}{\sinh t}=4\cosh^2t-1$ gives $(4\cdot\frac94-1)(4\cdot\frac{25}4-1)=8\cdot24=192$.

(c) $\tau_{m,n}=\tau_{n,m}$, so we may assume $m\le n$. Since $\frac{\sinh nt}{\sinh t}=e^{(n-1)t}\frac{1-e^{-2nt}}{1-e^{-2t}}$,
\[
\frac{\log\tau_{m,n}}{mn}=\frac{n-1}n\cdot\frac1m\sum_{j=1}^{m-1}t_j+E,\qquad|E|\le\frac2{mn}\sum_{j=1}^{m-1}\bigl|\log(1-e^{-2t_j})\bigr|.
\]
Now $\cosh t_j-1\ge1-\cos\frac\pi m\ge\frac2{m^2}$ and $\cosh t-1=2\sinh^2\frac t2\le\frac32t^2$ for $\cosh t\le3$, so $t_j\ge\frac1m$; and $1-e^{-2t}\ge\min(t,\frac12)$. Hence
$|E|\le\frac{2(\log m+1)}n\le\frac{2(\log n+1)}n\to0$. The main term is a Riemann sum of the continuous function $\theta\mapsto\operatorname{arccosh}(2-\cos\theta)$ on $[0,\pi]$, with step $\frac\pi m$.

(d) With $\theta=2u$, $2-\cos2u=1+2\sin^2u$, and $\operatorname{arccosh}(1+2s^2)=2\operatorname{arcsinh}s$ because $\cosh2y=1+2\sinh^2y$. So the limit is $\frac4\pi I(1)$, where
$I(\alpha)=\int_0^{\pi/2}\operatorname{arcsinh}(\alpha\sin u)\dd u$. Differentiating under the integral (the derivative is bounded by $1$),
\[
I'(\alpha)=\int_0^{\pi/2}\frac{\sin u\dd u}{\sqrt{1+\alpha^2\sin^2u}}=\int_0^1\frac{\dd c}{\sqrt{1+\alpha^2-\alpha^2c^2}}=\frac1\alpha\arcsin\frac\alpha{\sqrt{1+\alpha^2}}=\frac{\arctan\alpha}\alpha,
\]
with $c=\cos u$. So $I(1)=\int_0^1\frac{\arctan\alpha}\alpha\dd\alpha=\sum_{k\ge0}\frac{(-1)^k}{(2k+1)^2}=G$, integrating the alternating series $\sum_k\frac{(-1)^k\alpha^{2k}}{2k+1}$ term by term (its partial
sums are bounded by $1$). Finally $e^{4G/\pi}=3.20991\ldots$

\tools{Tome~11, Corollary~A.5 (matrix-tree theorem); Tome~2, Chapter~8 (Laplacians of products); Tome~6, Theorem~7.9 (dominated convergence).}
\leads The constant $\frac{4G}\pi$ is the \emph{entropy} of spanning trees of $\Z^2$. Kasteleyn and, independently, Temperley and Fisher counted domino tilings of the grid in 1961 with a product of the same
kind, and found the entropy $\frac G\pi$; Temperley's bijection (1974) between spanning trees and dominoes explains the coincidence of Catalan's constant. Burton and Pemantle (1993) described the
uniform random spanning tree of the plane.
\end{solution}

\begin{solution}{pb:squaredrum}
\idea The square drum separates into two strings, so its eigenvalues are sums of two squares, and counting the ways of writing $k=m^2+n^2$ is Jacobi's problem from Tome~1.

\writeup (a) The $\sqrt2\sin m\pi x$, $m\ge1$, form a Hilbert basis of $L^2(]0,1[)$ (Tome~7, Lemma~4.10), so the products are orthonormal by Fubini. Let $f\in L^2(\Omega)$ be orthogonal to all of
them, and $g_n(x)=\int_0^1f(x,y)\sqrt2\sin n\pi y\dd y$, which is in $L^2(]0,1[)$ by Cauchy--Schwarz. By Fubini, $g_n$ is orthogonal to every $\sqrt2\sin m\pi x$, so $g_n=0$ almost everywhere.
Outside a countable union of null sets, $f(x,\cdot)\in L^2$ and all its sine coefficients vanish, so $f(x,\cdot)=0$ almost everywhere, and $\iint|f|^2=0$ by Tonelli.

(b) An eigenfunction $u$ for $\lambda$ satisfies $\langle u,e\rangle(\lambda-\lambda_e)=0$ for each basis function $e$ with eigenvalue $\lambda_e$ (integrate by parts twice), so it lies in the span of the
products with $\lambda_e=\lambda$, and the multiplicity of $\pi^2k$ is the number of pairs $(m,n)$ with $m,n\ge1$ and $m^2+n^2=k$. Among the $r_2(k)=4\sum_{d\mid k}\chi_4(d)$ representations by
pairs in $\Z^2$, the four with a zero coordinate occur exactly when $k$ is a square, and the others split evenly among the four quadrants. So the multiplicity is
$\frac14\bigl(r_2(k)-4\cdot\mathbf 1_{k\text{ is a square}}\bigr)$.

(c) $s(k)=\sum_{d\mid k}\chi_4(d)$ is multiplicative, equal to $e+1$ at $p^e$ with $p\equiv1\pmod4$, to $1$ at $2^e$, and to $0$ or $1$ at $q^e$ with $q\equiv3\pmod4$ as $e$ is odd or even. At
$k=5^j$ the multiplicity is $j+1-\mathbf 1_{j\text{ even}}\ge j$. Multiplicity $3$ needs $s(k)=3$ with $k$ not a square, or $s(k)=4$ with $k$ a square; the second is impossible, since for a
square every exponent is even and $s(k)$ is a product of odd numbers. So $k=2^a p^2\cdot(\text{square})$ with $a$ odd, and the least is $2\cdot5^2=50=1+49=25+25$. Multiplicity $4$ needs
$s(k)=4$ with $k$ not a square, so $k$ has $p_1p_2$ or $p^3$ in it, or $s(k)=5$ with $k$ a square, so $p^4\mid k$; the least is $5\cdot13=65=1+64=16+49$.

\tools{Tome~7, Lemma~4.10 (the sine basis); Tome~1, Chapter~10 (Jacobi's two-squares formula, proved with Gaussian integers).}
\leads For a generic rectangle the multiplicities are all $1$; the square owes its large ones to arithmetic. The multiplicity of $\pi^2k$ is unbounded but small on average: $r_2(k)$ is $0$ for most
$k$ and $\pi$ on average (the Gauss circle problem, Tome~1, Chapter~7).
\end{solution}

\begin{solution}{pb:weylrect}
\idea Counting eigenvalues is counting lattice points in a quarter ellipse. Dirichlet counts the points off the axes, Neumann all of them; unit squares turn the counts into areas from both sides.

\writeup (a) $N_N(\lambda)$ counts $(m,n)\in\N^2$ in $E=\{(u,v)\in[0,\infty[^2:\frac{u^2}{a^2}+\frac{v^2}{b^2}\le\frac\lambda{\pi^2}\}$, and $N_D(\lambda)$ those with $m,n\ge1$. The points with $m=0$ are
$0\le n\le\frac{b\sqrt\lambda}\pi$, those with $n=0$ are $0\le m\le\frac{a\sqrt\lambda}\pi$, and $(0,0)$ is counted twice.

(b) $E$ is a quarter of an ellipse with semi-axes $\frac{a\sqrt\lambda}\pi$ and $\frac{b\sqrt\lambda}\pi$, of area $\frac\pi4\cdot\frac{ab\lambda}{\pi^2}=\frac{ab}{4\pi}\lambda$. If $(u,v)\in E$ and
$0\le u'\le u$, $0\le v'\le v$, then $(u',v')\in E$. So for $(m,n)\in E$ with $m,n\ge1$, the square $[m-1,m]\times[n-1,n]$ lies in $E$; these squares do not overlap, so $N_D\le\operatorname{area}E$.
Conversely, a point $(u,v)\in E$ lies in the square $[m,m+1]\times[n,n+1]$ with $m=\lfloor u\rfloor$, $n=\lfloor v\rfloor$, and $(m,n)\in E$: these squares cover $E$, so
$\operatorname{area}E\le N_N$.

(c) By (a) and (b), $0\le\frac{ab}{4\pi}\lambda-N_D(\lambda)\le N_N(\lambda)-N_D(\lambda)\le\frac{(a+b)\sqrt\lambda}\pi+1$. So $ab=\lim4\pi N_D(\lambda)/\lambda$, and the same limit with $N_N$. The perimeter is
$2(a+b)=2\pi\lim_{\lambda\to\infty}\frac{N_N(\lambda)-N_D(\lambda)}{\sqrt\lambda}$.

\tools{Elementary counting; the area of an ellipse.}
\leads Weyl proved in 1911--1912 that $N_D(\lambda)\sim\frac{\operatorname{area}(\Omega)}{4\pi}\lambda$ for every bounded domain, by squeezing it between unions of squares, a refinement of (b) with
Dirichlet and Neumann conditions on the small squares. He conjectured the second term $-\frac{\operatorname{length}(\partial\Omega)}{4\pi}\sqrt\lambda$; Ivrii proved it in 1980 under a condition
on periodic billiard trajectories. Figure~\ref{fig:weyl} shows it for one rectangle.
\end{solution}

\begin{solution}{pb:heatrect}
\idea The heat trace smooths the staircase $N_D$, and Poisson summation computes it to exponential accuracy: the area, the perimeter and the four corners each leave their own term.

\writeup (a) Poisson summation for the Gaussian (Tome~4, Chapter~8) gives $\theta(s)=\sum_{m\in\Z}e^{-\pi sm^2}=s^{-1/2}\theta(1/s)$. Hence
\[
\sum_{m\ge1}e^{-\pi sm^2}=\frac{\theta(s)-1}2=\frac1{2\sqrt s}-\frac12+\varepsilon(s),\qquad\varepsilon(s)=\frac1{\sqrt s}\sum_{m\ge1}e^{-\pi m^2/s},
\]
and $0\le\varepsilon(s)\le s^{-1/2}\sum_{m\ge1}e^{-\pi m/s}$, which is the bound since $m^2\ge m$.

(b) $Z_D(t)=S(s_a)S(s_b)$ with $S(s)=\sum_{m\ge1}e^{-\pi sm^2}$ and $s_a=\frac{\pi t}{a^2}$, $s_b=\frac{\pi t}{b^2}$, because $t\pi^2\frac{m^2}{a^2}=\pi s_am^2$. With $A=\frac a{2\sqrt{\pi t}}=\frac1{2\sqrt{s_a}}$
and $B$ likewise, (a) gives
\[
Z_D(t)=\bigl(A-\tfrac12+\varepsilon_a\bigr)\bigl(B-\tfrac12+\varepsilon_b\bigr)=AB-\frac{A+B}2+\frac14+R,
\]
where $R=\varepsilon_a(B-\frac12+\varepsilon_b)+\varepsilon_b(A-\frac12)$. For $t\le1$, $\varepsilon_a=O(t^{-1/2}e^{-a^2/t})$ since $e^{-\pi/s_a}=e^{-a^2/t}$, and $A,B=O(t^{-1/2})$, so
$R=O(t^{-1}e^{-c/t})$. Finally $AB=\frac{ab}{4\pi t}$ and $\frac{A+B}2=\frac{a+b}{4\sqrt{\pi t}}$.

(c) The spectrum determines $Z_D$, hence $ab=\lim_{t\to0}4\pi tZ_D(t)$ and then $a+b=\lim_{t\to0}4\sqrt{\pi t}\bigl(\frac{ab}{4\pi t}-Z_D(t)\bigr)$. The sides are the roots of $X^2-(a+b)X+ab$.
For Neumann, $\sum_{m\ge0}e^{-\pi sm^2}=\frac1{2\sqrt s}+\frac12+\varepsilon(s)$, so $Z_N(t)=\frac{ab}{4\pi t}+\frac{a+b}{4\sqrt{\pi t}}+\frac14+O(t^{-1}e^{-c/t})$, and it determines the sides too.

\tools{Tome~4, Chapter~8 (Poisson summation, the theta function).}
\leads Kac asked the question of the title in \emph{Can one hear the shape of a drum?} (Amer. Math. Monthly, 1966), and showed with this kind of expansion that area and perimeter can be heard.
For a polygon the constant term is a sum over the corners, $\frac{\pi^2-\alpha^2}{24\pi\alpha}$ for a corner of angle $\alpha$, which is $\frac1{16}$ for a right angle: the $\frac14$ of (b) is four corners.
In 1992 Gordon, Webb and Wolpert found two different polygons with the same Dirichlet spectrum, so the answer to Kac's question is no in general.
\end{solution}

\begin{solution}{pb:phasespace}
\idea The oscillator is a ladder: $a^\dagger$ climbs one rung and raises the energy by $2$. Counting rungs below $\lambda$ gives half the area of a disc, the semiclassical rule of one quantum state per
area $2\pi$ of phase space.

\writeup (a) For $u$ in the space, $a^\dagger au=(x-D)(xu+u')=x^2u+xu'-u-xu'-u''=Hu-u$, and in the same way $aa^\dagger u=Hu+u$. So $H=a^\dagger a+1$ and
$Ha^\dagger=a^\dagger aa^\dagger+a^\dagger=a^\dagger(H+1)+a^\dagger=a^\dagger(H+2)$. Since $ae^{-x^2/2}=0$, $H\varphi_0=\varphi_0$, and by induction
$H\varphi_{n+1}=a^\dagger(H+2)\varphi_n=(2n+3)\varphi_{n+1}$. Finally $a^\dagger(he^{-x^2/2})=(2xh-h')e^{-x^2/2}$, which raises the degree by exactly one.

(b) $H$ is symmetric on the space (two integrations by parts, the boundary terms vanish thanks to the Gaussian), so eigenfunctions for distinct eigenvalues are orthogonal. By (a), the span of
$\varphi_0,\dots,\varphi_n$ is the span of $x^ke^{-x^2/2}$, $k\le n$. Tome~7, Theorem~5.5 says that the span of the $x^ke^{-\pi x^2}$ is dense in $L^2(\R)$; the unitary map
$f\mapsto(2\pi)^{-1/4}f(x/\sqrt{2\pi})$ sends $x^ke^{-\pi x^2}$ to a multiple of $x^ke^{-x^2/2}$, so these are dense too, and the normalised $\varphi_n$ form a Hilbert basis.

(c) $N(\lambda)=\#\{n\ge0:2n+1\le\lambda\}=\lfloor\frac{\lambda+1}2\rfloor\sim\frac\lambda2$, and the disc of radius $\sqrt\lambda$ has area $\pi\lambda$. For the string the eigenvalues are $(\frac{\pi n}L)^2$, so
$N(\lambda)=\lfloor\frac{L\sqrt\lambda}\pi\rfloor$, while $\{(x,\xi):0\le x\le L,\ \xi^2\le\lambda\}$ has area $2L\sqrt\lambda$.

(d) With $x=\lambda^{1/4}y$, the area of $\{\xi^2+x^4\le\lambda\}$ is $\int2\sqrt{\lambda-x^4}\dd x=4\lambda^{3/4}\int_0^1\sqrt{1-y^4}\dd y$, and with $y^4=w$,
$\int_0^1\sqrt{1-y^4}\dd y=\frac14B(\frac14,\frac32)=\frac{\Gamma(\frac14)\Gamma(\frac32)}{4\Gamma(\frac74)}$. Using $\Gamma(\frac74)=\frac34\Gamma(\frac34)$ and $\Gamma(\frac14)\Gamma(\frac34)=\pi\sqrt2$, the rule
predicts
\[
N(\lambda)\sim c\lambda^{3/4},\qquad c=\frac{\Gamma(\frac14)^2}{3\pi\sqrt{2\pi}}=0.5564\ldots,
\]
so the $n$-th eigenvalue should grow like $(n/c)^{4/3}$.

\tools{Tome~7, Theorem~5.5 (Hermite functions); Tome~4, Chapter~5 (the Beta function).}
\leads The rule is the Weyl law for Schrödinger operators $-\Delta+V$, a theorem for confining potentials; its sharper form, the Bohr--Sommerfeld condition, replaces $n$ by $n+\frac12$ and is exact
for the oscillator. In quantum mechanics $a$ and $a^\dagger$ are the annihilation and creation operators, and the ladder of (a) reappears in every representation of $\mathfrak{sl}_2$.
\end{solution}

\begin{solution}{pb:nodal}
\idea Every eigenfunction for $5\pi^2$ is a positive function times a function whose zero set is a straight line in suitable coordinates, and a line cuts a square in two.

\writeup (a) The eigenvalue $5\pi^2$ comes from $(1,2)$ and $(2,1)$. By $\sin2\alpha=2\sin\alpha\cos\alpha$,
\[
u=\alpha\sin\pi x\sin2\pi y+\beta\sin2\pi x\sin\pi y=2\sin\pi x\sin\pi y\,(\alpha\cos\pi y+\beta\cos\pi x),
\]
and the first factor is positive on $\Omega$. The map $(x,y)\mapsto(w,z)=(\cos\pi x,\cos\pi y)$ is a homeomorphism of $\Omega$ onto $]-1,1[^2$, and it sends $\{u\ne0\}$ to the complement of the line
$\alpha z+\beta w=0$, which passes through the centre. The two open half-planes meet the open square in two non-empty convex sets, so there are exactly two nodal domains.

(b) $\sin m\pi x$ is non-zero exactly off the points $\frac jm$, with alternating signs on the intervals between them. So $\{u\ne0\}$ is the union of the $mn$ open rectangles of the grid, each connected,
and two of them are separated by a line where $u=0$.

(c) In units of $\pi^2$, the ordered list up to $20$ is $2;\ 5,5;\ 8;\ 10,10;\ 13,13;\ 17,17;\ 18;\ 20,20$. The first indices and the products are:
\begin{center}
\begin{tabular}{lcccccccc}
\toprule
eigenvalue $/\pi^2$&$2$&$5$&$8$&$10$&$13$&$17$&$18$&$20$\\
$(m,n)$&$(1,1)$&$(1,2)$&$(2,2)$&$(1,3)$&$(2,3)$&$(1,4)$&$(3,3)$&$(2,4)$\\
first index $k$&$1$&$2$&$4$&$5$&$7$&$9$&$11$&$12$\\
domains $mn$&$1$&$2$&$4$&$3$&$6$&$4$&$9$&$8$\\
\bottomrule
\end{tabular}
\end{center}
Courant's bound holds in every column, with equality only for $k=1$, $2$ and $4$.

\tools{Elementary topology of the square.}
\leads Courant proved his theorem in 1923, and Pleijel showed in 1956 that equality holds only for finitely many $k$ on every domain; for the square, $k=1,2,4$ are the only ones, a fact whose complete proof was written down by Bérard and Helffer. Pleijel's
proof combines Courant's argument with the Faber--Krahn inequality and Weyl's law.
\end{solution}

\begin{solution}{pb:eulerrational}
\idea Iterating the Green kernel of the string stays inside polynomials with rational coefficients, and its trace is a power sum of the eigenvalues $(n\pi)^{-2}$: rational numbers on one side,
$\zeta(2k)/\pi^{2k}$ on the other.

\writeup (a) $G_1$ is $x(1-y)$ on $x\le y$ and $y(1-x)$ on $x\ge y$. Suppose (a) holds for $G_k$, and fix $x\le y$. On each of $[0,x]$, $[x,y]$, $[y,1]$, both $G_k(x,z)$ and $G(z,y)$ are fixed polynomials
in $(x,y,z)$ with rational coefficients, and integrating such a polynomial in $z$ between two of the bounds $0,x,y,1$ gives a polynomial in $(x,y)$ with rational coefficients. The same for $x\ge y$.

(b) Let $T$ be the operator with kernel $G$; by Fubini $T^k$ has kernel $G_k$, and $G_k$ is symmetric because $T^k$ is self-adjoint. For operators with square-integrable kernels $a,b$, the
Hilbert--Schmidt inner product is $\iint a\overline b=\sum_n\langle Ae_n,Be_n\rangle$ (Tome~7, Proposition~4.13, polarised). With $A=T^{k-1}$, $B=T$ and $e_n=\sqrt2\sin n\pi x$,
\[
\int_0^1G_k(x,x)\dd x=\iint G_{k-1}(x,y)G(y,x)\dd y\dd x=\sum_n\langle T^{k-1}e_n,Te_n\rangle=\sum_{n\ge1}(n\pi)^{-2k}.
\]

(c) For $k\ge2$, $\zeta(2k)/\pi^{2k}=\int_0^1G_k(x,x)\dd x$, the integral of a polynomial with rational coefficients, hence rational. For $k=1$, $\zeta(2)=\frac{\pi^2}6$ (Tome~7, Exercise~4.6).

(d) For $x\le y$, $\int_0^1G(x,z)G(z,y)\dd z$ splits into
\[
\begin{aligned}
&\textstyle\int_0^xz(1-x)\,z(1-y)\dd z+\int_x^yx(1-z)\,z(1-y)\dd z+\int_y^1x(1-z)\,y(1-z)\dd z\\
&\qquad\textstyle=x(1-y)\Bigl[\frac{(1-x)x^2}3+\frac{y^2-x^2}2-\frac{y^3-x^3}3+\frac{y(1-y)^2}3\Bigr]=\frac16x(1-y)(2y-x^2-y^2).
\end{aligned}
\]
By (b) with $k=3$ and symmetry,
\[
\frac{\zeta(6)}{\pi^6}=2\iint_{x<y}G_2(x,y)\,x(1-y)=\frac13\int_0^1(1-y)^2\Bigl(\frac{2y^4}3-\frac{8y^5}{15}\Bigr)\dd y=\frac13\Bigl(\frac2{315}-\frac8{15\cdot168}\Bigr)=\frac1{945},
\]
using $\int_0^1y^4(1-y)^2=\frac1{105}$ and $\int_0^1y^5(1-y)^2=\frac1{168}$.

\tools{Tome~7, Proposition~4.13 (Hilbert--Schmidt operators), Lemma~4.10 and Exercise~4.6.}
\leads Euler found $\zeta(2k)=\frac{(-1)^{k+1}B_{2k}(2\pi)^{2k}}{2(2k)!}$ with the Bernoulli numbers (Tome~4). Nothing like this is known for odd arguments: Apéry proved in 1978 that $\zeta(3)$ is
irrational, and the traces of powers of $G$ only ever produce the even values.
\end{solution}

\begin{solution}{pb:eigendecay}
\idea Compressing to step functions loses little trace when the kernel is Lipschitz, and the trace lost bounds the tail of the eigenvalues. A small tail forces small eigenvalues.

\writeup (a) $PKP$ is positive self-adjoint on the range of $P$, with eigenvalues $\mu_1\ge\dots\ge\mu_m\ge0$, and $\langle PKPf,f\rangle=\langle Kf,f\rangle$ for $f$ in that range. By Courant--Fischer
(Tome~7, Theorem~4.6), $\mu_j$ is a maximum of minima of $\langle Kf,f\rangle$ over $j$-dimensional subspaces of the range of $P$, and $\lambda_j$ the same maximum over all $j$-dimensional
subspaces. So $\mu_j\le\lambda_j$ and $\operatorname{tr}(PKP)=\sum_j\mu_j\le\lambda_1+\dots+\lambda_m$.

(b) Let $I_i=[\frac{i-1}m,\frac im[$ and $e_i=\sqrt m\mathbf 1_{I_i}$, an orthonormal basis of the range. Then $\operatorname{tr}(PKP)=\sum_i\langle Ke_i,e_i\rangle=\sum_im\iint_{I_i^2}k(x,y)$, while by Mercer (Tome~7,
Theorem~4.14) $\sum_j\lambda_j=\int_0^1k(x,x)\dd x=\sum_im\iint_{I_i^2}k(x,x)\dd x\dd y$. So
\[
\sum_{j>m}\lambda_j\le\sum_j\lambda_j-\operatorname{tr}(PKP)=\sum_im\iint_{I_i^2}\bigl(k(x,x)-k(x,y)\bigr)\le\sum_imL\iint_{I_i^2}|x-y|=m\cdot mL\cdot\frac1{3m^3}=\frac L{3m},
\]
since $\iint_{[0,h]^2}|x-y|=\frac{h^3}3$.

(c) The eigenvalues decrease, so $m\lambda_{2m}\le\sum_{m<j\le2m}\lambda_j\le\frac L{3m}$ and $\lambda_{2m}\le\frac L{3m^2}$. For $n=2m$ this is $\frac{4L}{3n^2}$; for $n=2m+1$ with $m\ge1$,
$\lambda_n\le\lambda_{2m}\le\frac L{3m^2}\le\frac{3L}{n^2}$ since $n\le3m$. The kernel $\min(x,y)$ is positive, Lipschitz with $L=1$, and its eigenvalues are $\frac1{(n-\frac12)^2\pi^2}$ (Tome~7,
Example~4.15), so $n^2\lambda_n\to\frac1{\pi^2}$: no bound $o(n^{-2})$ holds.

\tools{Tome~7, Theorems~4.6 and~4.14, Example~4.15.}
\leads Smoother kernels have faster decaying eigenvalues, and analytic kernels have geometrically decaying ones. This is why a Gaussian process with a smooth covariance is nearly
finite-dimensional: by the Karhunen--Loève expansion it is a few Gaussian coordinates plus a tiny remainder, which low-rank methods based on singular-value truncation (Tome~2, Appendix~F), exploit.
\end{solution}

\begin{solution}{pb:slepian}
\idea $K_c$ first restricts a signal to the frequencies $|\xi|\le c$, then to the time interval $]-1,1[$. Its eigenvalues measure how well a signal can be concentrated in both. The trace counts
$\frac{2c}\pi$ degrees of freedom, and the Hilbert--Schmidt norm shows that almost all eigenvalues are near $0$ or $1$.

\writeup (a) $\frac1{2\pi}\int_{-c}^ce^{i\xi(x-y)}\dd\xi=\frac{\sin c(x-y)}{\pi(x-y)}$. For $f\in L^2(]-1,1[)$, extended by $0$, Fubini gives
$\langle K_cf,f\rangle=\frac1{2\pi}\int_{-c}^c|\hat f(\xi)|^2\dd\xi$ with $\hat f(\xi)=\int f(x)e^{-i\xi x}\dd x$. This is $\ge0$, and at most $\frac1{2\pi}\int_\R|\hat f|^2=\|f\|_2^2$ by Plancherel (Tome~7,
Theorem~5.4, in this normalisation). $K_c$ is compact and self-adjoint, so its eigenvalues are Rayleigh quotients, in $[0,1]$.

(b) The kernel is continuous and positive by (a), so Mercer's theorem gives $\sum_n\lambda_n=\int_{-1}^1\frac c\pi\dd x=\frac{2c}\pi$.

(c) $\sum_n\lambda_n^2=\iint K_c^2$. With $u=x-y$, whose density on $[-1,1]^2$ is $2-|u|$ on $[-2,2]$, and then $v=cu$,
\[
\sum_n\lambda_n^2=\frac2{\pi^2}\int_0^2(2-u)\frac{\sin^2cu}{u^2}\dd u=\frac2{\pi^2}\Bigl(2c\int_0^{2c}\frac{\sin^2v}{v^2}\dd v-\int_0^{2c}\frac{\sin^2v}v\dd v\Bigr).
\]
Now $\int_0^\infty\frac{\sin^2v}{v^2}=\frac\pi2$ and $\int_X^\infty\frac{\sin^2v}{v^2}=\frac1{2X}-\int_X^\infty\frac{\cos2v}{2v^2}=\frac1{2X}+O(X^{-2})$ (integrate by parts), so the first term in the bracket is
$\pi c-\frac12+O(\frac1c)$. The second is $\frac12\int_0^{4c}\frac{1-\cos w}w\dd w=\frac12(\gamma+\log4c)+o(1)$, by the classical formula $\int_0^Y\frac{1-\cos w}w\dd w=\gamma+\log Y-\operatorname{Ci}(Y)$
with $\operatorname{Ci}(Y)\to0$ (DLMF, §6.2). Hence $\sum\lambda_n^2=\frac{2c}\pi-\frac{1+\gamma+\log4c}{\pi^2}+o(1)$.

(d) By (b) and (c), $\sum_n\lambda_n(1-\lambda_n)=\frac{\log4c+\gamma+1}{\pi^2}+o(1)=O(\log c)$. Let $N=\#\{n:\lambda_n>\frac12\}$. Then
$N-\sum_n\lambda_n=\sum_{\lambda_n>1/2}(1-\lambda_n)-\sum_{\lambda_n\le1/2}\lambda_n$, and $1-\lambda\le2\lambda(1-\lambda)$ for $\lambda>\frac12$, $\lambda\le2\lambda(1-\lambda)$ for $\lambda\le\frac12$. So
$|N-\frac{2c}\pi|\le2\sum_n\lambda_n(1-\lambda_n)=O(\log c)$. On $[\varepsilon,1-\varepsilon]$, $\lambda(1-\lambda)\ge\varepsilon(1-\varepsilon)$, so at most $\frac{O(\log c)}{\varepsilon(1-\varepsilon)}$ eigenvalues lie there.

\tools{Tome~7, Theorems~4.14 (Mercer) and~5.4 (Plancherel).}
\leads Slepian, Pollak and Landau found in 1961--1962, at Bell Labs, that $K_c$ commutes with a differential operator whose eigenfunctions are the prolate spheroidal wave functions; this explains
the engineers' rule that a signal of bandwidth $W$ and duration $T$ has about $2WT$ degrees of freedom. Landau and Widom (1980) found the exact width of the plunge: the number of eigenvalues
in $]\varepsilon,1-\varepsilon[$ is $\frac2{\pi^2}\log\frac{1-\varepsilon}\varepsilon\log c+o(\log c)$.
\end{solution}

\begin{solution}{pb:jentzsch}
\idea A positive kernel prefers positive functions: replacing $f$ by $|f|$ can only raise the Rayleigh quotient. So the top eigenfunction can be taken positive, and then positivity of $k$ makes it
strictly positive, which forces simplicity and a gap to the rest of the spectrum.

\writeup (a) $K$ is compact and self-adjoint, and $\langle K\mathbf 1,\mathbf 1\rangle=\iint k>0$, so $\lambda_1>0$, and the largest eigenvalue of a compact self-adjoint operator, when positive, equals the
supremum of the Rayleigh quotient and is attained exactly on its eigenspace (expand in an eigenbasis; Tome~7, Theorem~4.6). If $Kf=\lambda_1f$ with $\|f\|=1$, then
$\langle K|f|,|f|\rangle=\iint k(x,y)|f(x)||f(y)|\ge|\langle Kf,f\rangle|=\lambda_1$, so $|f|$ attains the supremum and $K|f|=\lambda_1|f|$.

(b) $f=\lambda_1^{-1}Kf$ is continuous, as $k$ is uniformly continuous. By (a), $|f|=\lambda_1^{-1}K|f|$, and $(K|f|)(x)=\int k(x,y)|f(y)|\dd y>0$ for every $x$, since $k>0$ and $|f|$ is not almost
everywhere $0$. So $f$ never vanishes and, being continuous on $[0,1]$, has constant sign. If $f,g$ were independent eigenfunctions for $\lambda_1$, $g-\frac{g(0)}{f(0)}f$ would be one that vanishes at $0$.

(c) Let $Kg=\mu g$, $g\ne0$, $\mu\ne\lambda_1$. Then $|\mu|\|g\|^2=|\langle Kg,g\rangle|\le\langle K|g|,|g|\rangle\le\lambda_1\|g\|^2$. Suppose $\mu=-\lambda_1$. Equality holds throughout, so
$K|g|=\lambda_1|g|$ by (a), and for every $x$,
\[
\lambda_1|g(x)|=|(Kg)(x)|=\Bigl|\int k(x,y)g(y)\dd y\Bigr|\le\int k(x,y)|g(y)|\dd y=\lambda_1|g(x)|.
\]
Equality in the middle with $k>0$ forces $g\ge0$ almost everywhere or $g\le0$ almost everywhere; then $Kg=\pm K|g|=\lambda_1g$, so $\mu=\lambda_1$, a contradiction.

(d) Let $\varphi>0$ with $K\varphi=\lambda_1\varphi$, and $\alpha=\min\frac{Kg}g$. Then $Kg\ge\alpha g$ pointwise, and pairing with $\varphi>0$, $\lambda_1\langle g,\varphi\rangle=\langle g,K\varphi\rangle=\langle Kg,\varphi\rangle\ge\alpha\langle g,\varphi\rangle$,
with $\langle g,\varphi\rangle>0$. The upper bound is the same. Finally $\lambda_1\ge\langle K\mathbf 1,\mathbf 1\rangle=\iint k$.

\tools{Tome~7, Theorem~4.6; Tome~2, Chapter~7 (Perron--Frobenius and Collatz--Wielandt for matrices).}
\leads Jentzsch proved the theorem in 1912 (J. reine angew. Math. 141), five years after Perron's theorem for positive matrices and in the same year as Frobenius's work on non-negative ones.
Krein and Rutman (1948) extended it to operators preserving a cone in a Banach space. It is the reason why the multitype branching process of Problem~\ref{pb:multitype} and the transfer operators
of statistical mechanics have one dominant growth rate.
\end{solution}

\begin{solution}{pb:prediction}
\idea The prediction error is a ratio of consecutive Toeplitz determinants, so the determinant is a product of prediction errors. Szegő's theorem on the average of their logarithms then says where
the decreasing sequence $\sigma_n^2$ ends: at the geometric mean of the spectral density.

\writeup (a) By stationarity, the covariance matrix of $(X_0,X_{-1},\dots,X_{-n})$ is $T_{n+1}(f)=\begin{pmatrix}c_0&r^{\mathsf T}\\r&T_n(f)\end{pmatrix}$ with $r=(c_1,\dots,c_n)^{\mathsf T}$, and it is positive definite
by Problem~\ref{pb:szego}(a) since $\min f>0$. The error is $c_0-2a^{\mathsf T}r+a^{\mathsf T}T_n(f)a$, a positive definite quadratic function of $a$, minimal at $a=T_n(f)^{-1}r$ with value $c_0-r^{\mathsf T}T_n(f)^{-1}r$.
This is the Schur complement of $T_n(f)$ in $T_{n+1}(f)$, so $\det T_{n+1}(f)=\det T_n(f)\cdot\sigma_n^2$. The argument only uses that $T_{n+1}$ is positive definite.

(b) Taking $a_{n+1}=0$ shows $\sigma_{n+1}^2\le\sigma_n^2$. Also $\sigma_n^2=1/(T_{n+1}(f)^{-1})_{00}\ge\lambda_{\min}(T_{n+1}(f))\ge\min f>0$, so $\sigma_n^2\to\sigma^2>0$. By (a),
$\log\det T_n(f)=\sum_{j=0}^{n-1}\log\sigma_j^2$, a Cesàro sum, so $\frac1n\log\det T_n(f)\to\log\sigma^2$. By Problem~\ref{pb:szegodet}(a) the limit is $\int\log f$.

(c) For $\rho^{|k|}$, $\det A_n=(1-\rho^2)^{n-1}$ (Problem~\ref{pb:kms}), so $\sigma_n^2=1-\rho^2$ for every $n\ge1$, which is $\exp\int\log P_\rho$. The sequence is autoregressive: $X_0-\rho X_{-1}$ is
uncorrelated with the past, so the last value says everything and the error is reached at once. For $f=|1-e(\theta)|^2$, $\det T_n=n+1$ (Problem~\ref{pb:szegodet}(c)), so $\sigma_n^2=\frac{n+2}{n+1}\to1$.
Here $X_k=\xi_k-\xi_{k-1}$ for a white noise $\xi$ of variance $1$: the best one can do is to know $\xi_{-1}$, which leaves the error $1$, and a finite past recovers $\xi_{-1}$ only up to an error of
order $\frac1n$.

\tools{Problems~\ref{pb:szego}, \ref{pb:szegodet} and~\ref{pb:kms}; Tome~2, Chapter~11 (Schur complements); Tome~3, Theorem~6.15 (Gaussian vectors).}
\leads Kolmogorov (1941) and, independently, Wiener proved that for every stationary sequence with spectral density $f$ the error of prediction from the infinite past is $\exp\int\log f$; it is $0$,
and the future is a linear function of the past, exactly when $\int\log f=-\infty$. Wiener's version was a wartime report on anti-aircraft fire control, printed in 1949.
\end{solution}

\begin{solution}{pb:levinson}
\idea The new value $X_{-n-1}$ brings into the prediction of $X_0$ only its part that the values in between do not explain, the backward error. Projecting on it costs one inner product, and that
inner product, normalised, is the reflection coefficient.

\writeup Since the matrices are positive definite there is a centred Gaussian vector with these covariances (Tome~3, Theorem~6.15). Let $H_n$ be the span of $X_{-1},\dots,X_{-n}$, $e_f=X_0-\sum_ka^{(n)}_kX_{-k}$
the forward error and $e_b=X_{-n-1}-\sum_ka^{(n)}_kX_{-n-1+k}$. Reversing time preserves the covariances, since $\gamma$ is even, so $e_b$ is the backward prediction error: it is orthogonal to
$H_n$ and $\mathbb E e_b^2=\sigma_n^2$.

(a) $H_{n+1}=H_n\oplus\R e_b$, an orthogonal sum. So the prediction from $H_{n+1}$ is the prediction from $H_n$ plus the projection of $X_0$ on $e_b$, whose coefficient is
$\mathbb E[X_0e_b]/\sigma_n^2=\kappa_{n+1}$. Expanding $\kappa_{n+1}e_b$ gives the recursion for the coefficients. The new error is $e_f-\kappa_{n+1}e_b$, and since
$\mathbb E[e_fe_b]=\mathbb E[X_0e_b]=\kappa_{n+1}\sigma_n^2$, its variance is $\sigma_n^2-2\kappa_{n+1}^2\sigma_n^2+\kappa_{n+1}^2\sigma_n^2=(1-\kappa_{n+1}^2)\sigma_n^2$. It is positive, as in
Problem~\ref{pb:prediction}(a), so $|\kappa_{n+1}|<1$, and the product formula follows from $\sigma_0^2=\gamma(0)$.

(b) Suppose $\gamma(0),\dots,\gamma(n)$ are given, with $(\gamma(j-l))_{0\le j,l\le n}$ positive definite, realised by $X_{-n},\dots,X_0$, and let $\kappa\in\,]-1,1[$. Let $\xi\sim N(0,1)$ be independent, and put
\[
X_{-n-1}=\sum_{k=1}^na^{(n)}_kX_{-n-1+k}+\kappa e_f+\sqrt{1-\kappa^2}\,\sigma_n\xi.
\]
The first sum is the backward prediction, so it has the right covariances $\gamma(n+1-j)$ with $X_{-j}$, $1\le j\le n$, and $e_f$ and $\xi$ are orthogonal to $H_n$. The variance is
$(\gamma(0)-\sigma_n^2)+\kappa^2\sigma_n^2+(1-\kappa^2)\sigma_n^2=\gamma(0)$. And $\mathbb E[X_{-n-1}X_0]=\sum_ka^{(n)}_k\gamma(n+1-k)+\kappa\sigma_n^2$, which defines $\gamma(n+1)$ with $\kappa_{n+1}=\kappa$. The
new covariance matrix is positive definite, as $X_{-n-1}$ has an independent component of positive variance. Conversely, $\kappa_{n+1}$ is an affine function of $\gamma(n+1)$ with slope $\sigma_n^{-2}\ne0$,
so $\gamma(n+1)$ is determined by $\kappa_{n+1}$. Induction on $n$ gives existence and uniqueness.

(c) For $\rho^{|k|}$, $\kappa_1=\gamma(1)/\gamma(0)=\rho$, and $X_0-\rho X_{-1}$ is orthogonal to every $X_{-k}$, $k\ge1$ (as $\rho^k-\rho\cdot\rho^{k-1}=0$). So $a^{(n)}=(\rho,0,\dots,0)$ and
$\kappa_{n+1}=\sigma_n^{-2}(\rho^{n+1}-\rho\cdot\rho^n)=0$ for $n\ge1$. For the second sequence, let $a^{(n)}_k=-\frac{n+1-k}{n+1}$, a linear function of $k$ that equals $-1$ at $k=0$ and $0$ at $k=n+1$. The
normal equations $\sum_la_l\gamma(k-l)=\gamma(k)$ read $2a_k-a_{k-1}-a_{k+1}=\gamma(k)$ with $a_0,a_{n+1}$ omitted; the second difference of a linear function vanishes, so the left side is $a_0=-1$ for
$k=1$ and $a_{n+1}=0$ otherwise, as required. So $\kappa_n=a^{(n)}_n=-\frac1{n+1}$, and $\prod_{k\le n}(1-\frac1{(k+1)^2})=\frac{n+2}{2(n+1)}$ matches $\sigma_n^2=\frac{n+2}{n+1}$.

\tools{Tome~3, Theorem~6.15; Tome~7, Chapter~2 (orthogonal projections); Problem~\ref{pb:prediction}.}
\leads Levinson (1947) and Durbin (1960) found the recursion, which solves a Toeplitz system in $O(n^2)$ operations instead of $O(n^3)$; it is inside every speech codec that uses linear
prediction. (b) says that the $\kappa_n$, called Verblunsky or Schur coefficients, are free coordinates on the positive definite sequences. In the theory of orthogonal polynomials on the unit
circle they are the recursion coefficients, and Szegő's theorem becomes $\gamma(0)\prod_n(1-\kappa_n^2)=\exp\int\log f$.
\end{solution}

\begin{solution}{pb:quadratic}
\idea A Gaussian quadratic form is a weighted sum of independent $\chi^2$ variables, so its mean is a trace and its variance is twice a trace of a square. Szegő's theorem computes the first, and
operator norms bound the second.

\writeup (a) $\mathbb E[X^{\mathsf T}AX]=\sum_{j,l}A_{jl}\mathbb E[X_jX_l]=\operatorname{tr}(AT_n(f))$. By Problem~\ref{pb:szego}(b), $T_n(g)T_n(f)-T_n(gf)$ has at most $2d$ non-zero rows, $d$ the degree of $g$,
and by Problem~\ref{pb:szego}(a) its norm is at most $2\|f\|_\infty\|g\|_\infty$; so its trace is $O(1)$, while $\operatorname{tr}T_n(gf)=n\int fg$.

(b) Write $\Sigma=T_n(f)$, $A=T_n(g)$, $X=\Sigma^{1/2}Z$ with $Z$ standard, and $B=\Sigma^{1/2}A\Sigma^{1/2}=O\operatorname{diag}(\beta_i)O^{\mathsf T}$. Then $W=O^{\mathsf T}Z$ is standard (Tome~3, Theorem~6.15),
$X^{\mathsf T}AX=\sum_i\beta_iW_i^2$, and $\operatorname{Var}=\sum_i\beta_i^2\operatorname{Var}(W_i^2)=2\operatorname{tr}B^2=2\operatorname{tr}(A\Sigma A\Sigma)$. Finally
$\operatorname{tr}B^2\le n\|B\|^2\le n\|A\|^2\|\Sigma\|^2\le n\|g\|_\infty^2\|f\|_\infty^2$, and dividing by $n^2$ gives the bound.

(c) For $k\ge1$ take $g(\theta)=2\cos2\pi k\theta$, whose Toeplitz matrix has ones where $|j-l|=k$: then $\frac1nX^{\mathsf T}T_n(g)X=\frac2n\sum_{j=1}^{n-k}X_jX_{j+k}$, and $\int fg=2c_k$. By (a) and (b),
$\frac1n\sum_jX_jX_{j+k}$ has mean $c_k+O(\frac1n)$ and variance $O(\frac1n)$, so it tends to $c_k$ in $L^2$. For $k=0$ take $g=1$.

\tools{Problem~\ref{pb:szego}; Tome~3, Theorem~6.15.}
\leads The same computation, with $g=\frac1f$, is behind Whittle's approximation (1953) of the Gaussian likelihood of a stationary series by an integral of the periodogram, and Grenander and
Szegő's \emph{Toeplitz Forms and their Applications} (1958) is largely about such statistical uses.
\end{solution}

\begin{solution}{pb:ruin}
\idea The walk killed at $0$ and $N$ is a Toeplitz matrix of size $N-1$, diagonalised by sines. Survival is a sum of exponentials in $n$, and the slowest one, the lowest mode of the string, wins.

\writeup (a) Let $Q$ be the $(N-1)\times(N-1)$ matrix with $Q_{xy}=\frac12$ if $|x-y|=1$ and $0$ otherwise. A path of $n$ steps that stays in $\{1,\dots,N-1\}$ has probability given by a product of
entries of $Q$, so $P_x(T>n)=(Q^n\mathbf 1)(x)$. $Q=T_{N-1}(\cos2\pi\theta)$, so by Problem~\ref{pb:tridiag} its eigenvalues are $\cos\frac{\pi k}N$, with eigenvectors
$v_k(j)=\sqrt{2/N}\sin\frac{\pi kj}N$, orthonormal because $\sum_{j=1}^{N-1}\sin^2\frac{\pi kj}N=\frac{N-1}2-\frac12\sum_{j=1}^{N-1}\cos\frac{2\pi kj}N=\frac N2$. Expanding $\mathbf 1$ in this basis,
$Q^n\mathbf 1=\sum_k\cos^n\frac{\pi k}N\langle\mathbf 1,v_k\rangle v_k$, which is the formula with $c_k(x)=v_k(x)\langle\mathbf 1,v_k\rangle$.

(b) For $N=2$, $T=1$; let $N\ge3$. Since $\sin\frac{\pi(N-k)j}N=(-1)^{j+1}\sin\frac{\pi kj}N$, $v_{N-k}(j)=(-1)^{j+1}v_k(j)$, and $\cos^{2n}$ takes the same value at $k$ and $N-k$. So
$c_k(x)+c_{N-k}(x)=v_k(x)\sum_jv_k(j)\bigl(1+(-1)^{x+j}\bigr)=2v_k(x)\sum_{j\equiv x\,(2)}v_k(j)$. For $k=1$ every $v_1(j)$ is positive, so $C_x=2v_1(x)\sum_{j\equiv x\,(2)}v_1(j)>0$ (for $N=3$ the pair is $k=1,2$ and nothing else remains). All other $k$ have $|\cos\frac{\pi k}N|<\cos\frac\pi N$, so their contribution is $o(\cos^{2n}\frac\pi N)$. Finally
$\log\cos\frac\pi N=-\frac{\pi^2}{2N^2}+O(N^{-4})$, so $\cos^m\frac\pi N\approx e^{-\pi^2m/(2N^2)}$.

(c) $\|Q\|=\cos\frac\pi N<1$, so $\mathbb E_xT=\sum_{n\ge0}P_x(T>n)=((I-Q)^{-1}\mathbf 1)(x)$, which by (a) is $\sum_k\frac{c_k(x)}{1-\cos(\pi k/N)}$. And $h(x)=x(N-x)$ solves $(I-Q)h=\mathbf 1$:
$h(x)-\frac12(h(x-1)+h(x+1))=1$ for $1\le x\le N-1$, with $h(0)=h(N)=0$, which is exactly what the missing entries of $Q$ at the ends require. So $(I-Q)^{-1}\mathbf 1=h$.

\tools{Problem~\ref{pb:tridiag}; Tome~3, Chapter~8 (Markov chains).}
\leads The rate $\frac{\pi^2}{2N^2}$ is the first Dirichlet eigenvalue of $-\frac12\frac{\mathrm d^2}{\mathrm dx^2}$ on $[0,N]$: on large scales the walk is a Brownian motion, and the probability that
Brownian motion stays in an interval for a long time is governed by the lowest note of a string. The formula of (a) is in Feller's \emph{Introduction to Probability Theory}, volume~I, Chapter~XIV.
\end{solution}

\begin{solution}{pb:randomcirculant}
\idea A circulant is a polynomial in the cyclic shift, so the discrete Fourier transform diagonalises it, and its eigenvalues are random trigonometric sums. Parseval gives a floor, Hoeffding and a
union bound a ceiling.

\writeup (a) Let $S$ be the cyclic shift, $(Sx)_l=x_{l+1}$, indices modulo $n$. Then $C=\sum_j\varepsilon_jS^j$. The vectors $f_k(l)=n^{-1/2}e(\frac{kl}n)$ form an orthonormal basis with $Sf_k=e(\frac kn)f_k$,
so $Cf_k=\lambda_kf_k$. $C$ is diagonal in an orthonormal basis, hence normal, and the norm of a normal matrix is the largest modulus of an eigenvalue.

(b) $\sum_k|\lambda_k|^2=\operatorname{tr}(C^*C)=\sum_{j,l}|C_{jl}|^2=n^2$. So $\max_k|\lambda_k|^2\ge n$.

(c) $\operatorname{Re}\lambda_k=\sum_j\varepsilon_j\cos\frac{2\pi jk}n$ with $\sum_j\cos^2\le n$. Hoeffding's inequality (Tome~3, Theorem~4.9) gives $P(|\operatorname{Re}\lambda_k|\ge s)\le2e^{-s^2/(2n)}$, and the
same for the imaginary part. If $|\lambda_k|\ge t$, one of them is at least $\frac t{\sqrt2}$, so $P(|\lambda_k|\ge t)\le4e^{-t^2/(4n)}$, and the union bound over the $n$ values of $k$ gives the result.
For $t^2=8n\log n$, $4ne^{-2\log n}=\frac4n$.

\tools{Tome~4, Chapter~7 (the discrete Fourier transform); Tome~3, Theorem~4.9 (Hoeffding).}
\leads So $\sqrt n\le\|C\|\le\sqrt{8n\log n}$ with high probability. The truth is the upper order: $\|C\|$ is of order $\sqrt{n\log n}$ (Meckes, 2009), because $n$ nearly independent Gaussian-like
eigenvalues have a maximum of that size. For a random symmetric matrix with independent signs the norm is close to $2\sqrt n$: the circulant pays a logarithm for its rigid structure.
\end{solution}

\begin{solution}{pb:multitype}
\idea The mean matrix plays the role of the mean offspring number, and its Perron eigenvector turns the population vector into one number that grows exactly by the factor $\rho$ in expectation.

\writeup (a) Given $Z_n$, $Z_{n+1}$ is the sum over the $Z_n(i)$ individuals of each type $i$ of independent child vectors with mean $M_{i\cdot}$, so $\mathbb E[Z_{n+1}\mid Z_n]=Z_nM$ and
$\mathbb EZ_n=Z_0M^n$. Then $P(Z_n\ne0)\le\mathbb E|Z_n|_1\le|Z_0|_1\|M^n\|$ for a suitable matrix norm, and $\|M^n\|^{1/n}\to\rho<1$ by Gelfand's formula (Tome~7, Theorem~3.6). Since $0$ is absorbing, the
events $\{Z_n\ne0\}$ decrease, and their intersection has probability $0$.

(b) $\mathbb E[Z_{n+1}v\mid Z_n]=Z_nMv=\rho Z_nv$, so $W_n$ is a martingale. Given $Z_n$, $Z_{n+1}v$ is a sum of $|Z_n|_1$ independent terms, each with variance at most $s^2$, the largest of the
variances of $\xi\cdot v$ over the types. Hence
\[
\mathbb E(W_{n+1}-W_n)^2=\rho^{-2n-2}\,\mathbb E\operatorname{Var}(Z_{n+1}v\mid Z_n)\le\rho^{-2n-2}s^2\,\mathbb E|Z_n|_1\le\rho^{-2n-2}s^2\frac{Z_0M^nv}{\min_iv_i}=\frac{s^2Z_0v}{\rho^2\min_iv_i}\rho^{-n}.
\]
Martingale increments are orthogonal, so $\mathbb EW_n^2=W_0^2+\sum_{k<n}\mathbb E(W_{k+1}-W_k)^2$ stays bounded, as $\rho>1$.

(c) By Tome~3, Theorem~10.18, $W_n\to W$ almost surely and in $L^2$, so $\mathbb EW=W_0=Z_0v>0$ and $P(W>0)>0$. On $\{W>0\}$, $Z_nv\sim W\rho^n\to\infty$, so the population never dies, and
$|Z_n|_1$ lies between $\frac{Z_nv}{\max v_i}$ and $\frac{Z_nv}{\min v_i}$, both of order $\rho^n$.

\tools{Tome~7, Theorem~3.6 (Gelfand); Tome~2, Chapter~7 (Perron--Frobenius); Tome~3, Theorem~10.18 ($L^2$ martingales).}
\leads Harris's \emph{The Theory of Branching Processes} (1963) develops the multitype theory. Kesten and Stigum (1966) showed that $W>0$ almost surely on survival if and only if
$\mathbb E[|\xi|_1\log|\xi|_1]<\infty$ for every type, and that on survival $Z_n/|Z_n|_1$ tends to the normalised left Perron eigenvector: the proportions of types settle, whatever the start.
\end{solution}

\section*{Map of the chapter}

\begin{center}
\small
\adjustbox{max width=\linewidth}{\begin{tabular}{lll}
Problem & Topic & Tomes used\\\hline
\ref{pb:tridiag} & eigenvalues of tridiagonal Toeplitz matrices & 0, S1\\
\ref{pb:szego} & Szegő's first limit theorem & 6, 7, S1\\
\ref{pb:szegodet} & Toeplitz determinants and the strong theorem & 6\\
\ref{pb:kms} & the Kac--Murdock--Szegő matrix & 0\\
\ref{pb:redheffer} & Redheffer's matrix and the Mertens function & 0, 1\\
\ref{pb:minmatrix} & the matrix $\min(i,j)$ and $\sum(2k-1)^{-2}$ & 0, 6\\
\ref{pb:gridtrees} & spanning trees of the grid and Catalan's constant & C1, S1, 6\\
\ref{pb:squaredrum} & the square drum and sums of two squares & 1, 7\\
\ref{pb:weylrect} & Weyl's law for a rectangle & --\\
\ref{pb:heatrect} & hearing the shape of a rectangle & 4\\
\ref{pb:phasespace} & the harmonic oscillator and the phase-space rule & 4, 7\\
\ref{pb:nodal} & nodal domains of the square & --\\
\ref{pb:eulerrational} & Euler's theorem on $\zeta(2k)$ from a string & 7\\
\ref{pb:eigendecay} & eigenvalue decay of Lipschitz kernels & 7\\
\ref{pb:slepian} & time--frequency concentration & 7\\
\ref{pb:jentzsch} & Jentzsch's theorem & 7, S1\\
\ref{pb:prediction} & the Kolmogorov--Szegő formula & 3, S1\\
\ref{pb:levinson} & Levinson--Durbin and reflection coefficients & 3, 7\\
\ref{pb:quadratic} & estimating covariances & 3\\
\ref{pb:ruin} & the duration of the gambler's ruin & 3\\
\ref{pb:randomcirculant} & the norm of a random circulant & 3, 4\\
\ref{pb:multitype} & multitype branching & 3, 7, S1\\
\end{tabular}}
\end{center}

\section*{Sources and further reading}

\begin{itemize}
\item J.~J.~O'Connor and E.~F.~Robertson, Gábor Szegő, \emph{MacTutor History of Mathematics Archive}, University of St Andrews (Szegő's life, Pólya's words on the conjecture and on the book, and
Stuart Jay Sidney's review quoted in the epigraph); P.~Deift, A.~Its and I.~Krasovsky, Toeplitz matrices and Toeplitz determinants under the impetus of the Ising model: some history and some
recent results, \emph{Communications on Pure and Applied Mathematics} 66 (2013), 1360--1438, arXiv:1207.4990 (Onsager, Kaufman, Kakutani, Yang).
\item G.~Szegő, Ein Grenzwertsatz über die Toeplitzschen Determinanten einer reellen positiven Funktion, \emph{Mathematische Annalen} 76 (1915), 490--503; Beiträge zur Theorie der Toeplitzschen
Formen, \emph{Mathematische Zeitschrift} 6 (1920), 167--202; On certain Hermitian forms associated with the Fourier series of a positive function, \emph{Communications du Séminaire Mathématique de
l'Université de Lund}, tome supplémentaire (1952), 228--238 (Problems~\ref{pb:szego} and~\ref{pb:szegodet}).
\item G.~Pólya and G.~Szegő, \emph{Aufgaben und Lehrsätze aus der Analysis}, 2 vols., Springer, 1925; U.~Grenander and G.~Szegő, \emph{Toeplitz Forms and their Applications}, University of
California Press, 1958; A.~Böttcher and B.~Silbermann, \emph{Introduction to Large Truncated Toeplitz Matrices}, Springer, 1999.
\item M.~Kac, W.~L.~Murdock and G.~Szegő, On the eigenvalues of certain Hermitian forms, \emph{Journal of Rational Mechanics and Analysis} 2 (1953), 767--800 (Problem~\ref{pb:kms}).
\item R.~Redheffer, Eine explizit lösbare Optimierungsaufgabe, \emph{Internationale Schriftenreihe zur Numerischen Mathematik} 36 (1977), 213--216; W.~W.~Barrett, R.~W.~Forcade and A.~D.~Pollington,
On the spectral radius of a $(0,1)$ matrix related to Mertens' function, \emph{Linear Algebra and its Applications} 107 (1988), 151--159; R.~C.~Vaughan, On the eigenvalues of Redheffer's
matrix~I, in \emph{Number Theory with an Emphasis on the Markoff Spectrum}, Dekker, 1993 (Problem~\ref{pb:redheffer}).
\item P.~W.~Kasteleyn, The statistics of dimers on a lattice, \emph{Physica} 27 (1961), 1209--1225; H.~N.~V.~Temperley and M.~E.~Fisher, Dimer problem in statistical mechanics: an exact result,
\emph{Philosophical Magazine} 6 (1961), 1061--1063; R.~Burton and R.~Pemantle, Local characteristics, entropy and limit theorems for spanning trees and domino tilings via transfer-impedances,
\emph{Annals of Probability} 21 (1993), 1329--1371 (Problem~\ref{pb:gridtrees}).
\item H.~Weyl, Das asymptotische Verteilungsgesetz der Eigenwerte linearer partieller Differentialgleichungen, \emph{Mathematische Annalen} 71 (1912), 441--479; M.~Kac, Can one hear the shape of a
drum?, \emph{American Mathematical Monthly} 73 (1966), 1--23; C.~Gordon, D.~L.~Webb and S.~Wolpert, One cannot hear the shape of a drum, \emph{Bulletin of the AMS} 27 (1992), 134--138
(Problems~\ref{pb:weylrect} and~\ref{pb:heatrect}).
\item R.~Courant and D.~Hilbert, \emph{Methods of Mathematical Physics}, vol.~I, Interscience, 1953, Chapter~VI (nodal domains, the min-max principle); Å.~Pleijel, Remarks on Courant's nodal line
theorem, \emph{Communications on Pure and Applied Mathematics} 9 (1956), 543--550; P.~Bérard and B.~Helffer, Dirichlet eigenfunctions of the square membrane: Courant's property, and A.~Stern's
and Å.~Pleijel's analyses, in \emph{Analysis and Operator Theory}, Springer, 2019 (Problem~\ref{pb:nodal}).
\item D.~Slepian and H.~O.~Pollak, Prolate spheroidal wave functions, Fourier analysis and uncertainty~I, \emph{Bell System Technical Journal} 40 (1961), 43--63; H.~J.~Landau and H.~O.~Pollak, ibid.~II
(1961), 65--84, and~III, 41 (1962), 1295--1336; H.~J.~Landau and H.~Widom, Eigenvalue distribution of time and frequency limiting, \emph{Journal of Mathematical Analysis and Applications} 77 (1980),
469--481 (Problem~\ref{pb:slepian}).
\item R.~Jentzsch, Über Integralgleichungen mit positivem Kern, \emph{Journal für die reine und angewandte Mathematik} 141 (1912), 235--244; M.~G.~Krein and M.~A.~Rutman, Linear operators leaving
invariant a cone in a Banach space, \emph{Uspekhi Matematicheskikh Nauk} 3 (1948), 3--95 (Problem~\ref{pb:jentzsch}).
\item A.~N.~Kolmogorov, Stationary sequences in Hilbert space, \emph{Bulletin of Moscow State University, Mathematics} 2 (1941); N.~Wiener, \emph{Extrapolation, Interpolation, and Smoothing of
Stationary Time Series}, MIT Press, 1949; N.~Levinson, The Wiener RMS error criterion in filter design and prediction, \emph{Journal of Mathematics and Physics} 25 (1947), 261--278; J.~Durbin, The
fitting of time-series models, \emph{Revue de l'Institut International de Statistique} 28 (1960), 233--244; B.~Simon, \emph{Orthogonal Polynomials on the Unit Circle}, AMS, 2005
(Problems~\ref{pb:prediction} and~\ref{pb:levinson}).
\item P.~Whittle, Estimation and information in stationary time series, \emph{Arkiv för Matematik} 2 (1953), 423--434 (Problem~\ref{pb:quadratic}); W.~Feller, \emph{An Introduction to Probability Theory
and its Applications}, vol.~I, 3rd ed., Wiley, 1968, Chapter~XIV (Problem~\ref{pb:ruin}).
\item M.~W.~Meckes, Some results on random circulant matrices, in \emph{High Dimensional Probability~V}, IMS Collections 5 (2009), 213--223 (Problem~\ref{pb:randomcirculant}); T.~E.~Harris,
\emph{The Theory of Branching Processes}, Springer, 1963; H.~Kesten and B.~P.~Stigum, A limit theorem for multidimensional Galton--Watson processes, \emph{Annals of Mathematical Statistics} 37
(1966), 1211--1223 (Problem~\ref{pb:multitype}).
\end{itemize}
